% !TEX root = ../Elements.tex
% Greek and English text last checked 17/03/2013
%%%%%%
% BOOK 5
%%%%%%
\pdfbookmark[0]{Book 5}{book5}
\pagestyle{empty}
\begin{center}
{\Huge ELEMENTS BOOK 5}\\
\spa\spa\spa
{\huge\it Proportion\symbolfootnote[2]{The theory of proportion set out in this
book is generally attributed to Eudoxus of Cnidus. The novel feature of this theory is its ability to deal with irrational magnitudes, which had hitherto been
a major stumbling block for Greek mathematicians. Throughout the footnotes in this
book, $\alpha$, $\beta$, $\gamma$, {\rm etc.}, denote general (possibly irrational) magnitudes, whereas $m$, $n$, $l$,  {\rm etc.}, denote
positive integers.}}
\end{center}
%\vspace{3cm}
%\epsfysize=4in
%\centerline{\epsffile{Book05/front.eps}}
\newpage

%%%%%%%
% Definitions
%%%%%%%
\pdfbookmark[1]{Definitions}{def5}
\pagestyle{fancy}
\cfoot{{\greekfont τ̀ηεπαγε}}
\lhead{\large{\greekfont ΣΤΟΙΘΕΙΩΝ γ̀γν{5}.}}
\rhead{\large ELEMENTS BOOK 5}
\begin{Parallel}{}{} 
\ParallelLText{
\begin{center}
\large{{\greekfont Ὅροι}.}
\end{center}\vspace*{-7pt}

\ggn{1}.~{\greekfont Μέρος ἐστὶ μέγεθος μεγέθους τὸ ἔλασσον τοῦ μείζονος,
ὅταν καταμετρῇ τὸ μεῖζον.}

\ggn{2}.~{\greekfont Πολλαπλάσιον δὲ τὸ μεῖζον τοῦ ἐλάττονος, ὅταν καταμετρῆται ὑπὸ τοῦ ἐλάττονος.}

\ggn{3}.~{\greekfont Λόγος ἐστὶ δύο μεγεθῶν ὁμογενῶν ἡ κατὰ πηλικότητά
ποια σθέσις.}

\ggn{4}.~{\greekfont Λόγον ἔχειν πρὸς ἄλληλα μεγέθη λέγεται, ἃ δύναται
πολλαπλασιαζόμενα ἀλλήλων ὑπερέθειν.}

\ggn{5}.~{\greekfont Ἐν τῷ α ὐτῷ λόγῳ μεγέθη λέγεται εἶναι πρῶτον πρὸς
δεύτερον καὶ τρίτον πρὸς τέταρτον, ὅταν τὰ τοῦ πρώτου καί
τρίτου ἰσάκις πολλαπλάσια τῶν τοῦ δευτέρου καὶ τετάρτου
ἰσάκις πολλαπλασίων καθ ὁποιονοῦν πολλαπλασιασμὸν
ἑκάτερον ἑκατέρου ἢ ἅμα ὑπερέθῃ ἢ ἅμα ἴσα ᾖ ἢ
ἅμα ἐλλείπῇ ληφθέντα κατάλληλα.}

\ggn{6}.~{\greekfont Τὰ δὲ τὸν α ὐτὸν ἔθοντα λόγον μεγέθη ἀνάλογον
καλείσθω.}

\ggn{7}.~{\greekfont Ὅταν δὲ τῶν ἰσάκις πολλαπλασίων τὸ μὲν τοῦ πρώτου
πολλαπλάσιον ὑπερέθῃ τοῦ τοῦ δευτέρου πολλαπλασίου, τὸ δὲ
τοῦ τρίτου πολλαπλάσιον μὴ ὑπερέθῃ τοῦ τοῦ τετάρτου
πολλαπλασίου, τότε τὸ πρῶτον πρὸς τὸ δεύτερον μείζονα λόγον ἔχειν
λέγεται, ἤπερ τὸ τρίτον πρὸς τὸ τέταρτον.}

\ggn{8}.~{\greekfont Ἀναλογία δὲ ἐν τρισὶν ὅροις ἐλαθίστη ἐστίν.}

\ggn{9}.~{\greekfont Ὅταν δὲ τρία μεγέθη ἀνάλογον ᾖ, τὸ πρῶτον πρὸς τὸ
τρίτον διπλασίονα λόγον ἔχειν λέγεται ἤπερ πρὸς τὸ
δεύτερον.}

\ggn{10}.~{\greekfont Ὅταν δὲ τέσσαρα μεγέθη ἀνάλογον ᾖ, τὸ πρῶτον
πρὸς τὸ τέταρτον τριπλασίονα λόγον ἔχειν λέγεται ἤπερ
πρὸς τὸ δεύτερον, καὶ ἀεὶ ἑχῆς ὁμοίως, ὡς
ἂν ἡ ἀναλογία ὑπάρθῃ.}

\ggn{11}.~{\greekfont Ὁμόλογα μεγέθη λέγεται τὰ μὲν ἡγούμενα τοῖς
ἡγουμένοις τὰ δὲ ἑπόμενα τοῖς ἑπομένοις.}

\ggn{12}.~{\greekfont Ἐναλλὰχ λόγος ἐστὶ λῆψις τοῦ ἡγουμένου πρὸς τὸ
ἡγούμενον καὶ τοῦ ἑπομένου πρὸς τὸ ἑπόμενον.}

\ggn{13}.~{\greekfont Ἀνάπαλιν λόγος ἐστὶ λῆψις τοῦ ἑπομένου ὡς
ἡγουμένου πρὸς τὸ ἡγούμενον ὡς ἑπόμενον.}

\ggn{14}.~{\greekfont Σύνθεσις λόγου ἐστὶ λῆψις τοῦ ἡγουμένου μετὰ τοῦ
ἑπομένου ὡς ἑνὸς πρὸς α ὐτὸ τὸ ἑπόμενον.}

\ggn{15}.~{\greekfont Διαίρεσις λόγου ἐστὶ λῆψις τῆς ὑπεροθῆς, ᾗ ὑπερέθει
τὸ ἡγούμενον τοῦ ἑπομένου, πρὸς α ὐτὸ τὸ ἑπόμενον.}

\ggn{16}.~{\greekfont Ἀναστροφὴ λόγου ἐστὶ λῆψις τοῦ ἡγουμένου πρὸς
τὴν ὑπεροθήν, ῇ̔ ὑπερέθει τὸ ἡγούμενον τοῦ
ἑπομένου.}

\ggn{17}.~{\greekfont Δι ἴσου λόγος ἐστὶ πλειόνων ὄντων μεγεθῶν καὶ
ἄλλων α ὐτοῖς ἴσων τὸ πλῆθος σύνδυο λαμβανομένων
καὶ ἐν τῷ α ὐτῷ λόγῳ, ὅταν ᾖ ὡς ἐν τοῖς πρώτοις
μεγέθεσι τὸ πρῶτον πρὸς τὸ ἔσθατον, οὕτως ἐν τοῖς
δευτέροις μεγέθεσι τὸ πρῶτον πρὸς τὸ ἔσθατον; ἢ ἄλλως;
λῆψις τῶν ἄκρων καθ ὑπεχαίρεσιν τῶν  μέσων.}

\ggn{18}.~{\greekfont Τεταραγμένη δὲ ἀναλογία ἐστίν, ὅταν τριῶν ὄντων
μεγεθῶν καὶ ἄλλων α ὐτοῖς ἴσων τὸ πλῆθος
γίνηται ὡς μὲν ἐν τοῖς πρώτοις μεγέθεσιν ἡγούμενον πρὸς
ἐπόμενον, οὕτως  ἐν τοῖς δευτέροις μεγέθεσιν
ἡγούμενον πρὸς ἑπόμενον, ὡς δὲ ἐν τοῖς πρώτοις
μεγέθεσιν ἑπόμενον πρὸς ἄλλο τι, οὕτως ἐν τοῖς
δευτέροις ἄλλο τι πρὸς ἡγούμενον.}}

\ParallelRText{
\begin{center}
{\large Definitions}
\end{center}

1.~A magnitude is a part of a(nother) magnitude, the lesser of the greater, when
it measures the greater.$^\dag$

2.~And the greater (magnitude is) a multiple of the lesser when it is measured by
the lesser.

3.~A ratio is a certain type of relation  with respect to size of two magnitudes of the same kind.$^\ddag$

4.~(Those) magnitudes are said to have a ratio with respect to one another which,
being multiplied, are capable of exceeding one another.$^\S$

5.~Magnitudes are said to be in the same ratio, the first to the second, and
the third to the fourth, when  equal multiples of the first and the third 
either both exceed, are both equal to, or are both less than, 
equal multiples of the second and the fourth, respectively,  being taken in corresponding order, according to any kind of multiplication whatever.$^\P$

6.~And let magnitudes having the same ratio be called proportional.$^\ast$

7.~And when for equal multiples (as in Def.~5), the multiple of the first (magnitude) exceeds 
the multiple of the second, and the multiple of the third (magnitude) does not
exceed the multiple of the fourth, (then) the first (magnitude) is said to have a greater ratio
to the second than the third (magnitude has) to the fourth.

8.~And a  proportion in three terms is the smallest (possible).$^\$$

9.~And when three magnitudes are proportional, the first is said to
have to the third the squared$^\flat$ ratio  of that (it has) to the second.$^{\dag\dag}$

10.~And when four magnitudes are (continuously) proportional, the first is said to have to the fourth the cubed$^{\ddag\ddag}$ ratio  of that (it has) to the second.$^{\S\S}$ And  so on,   similarly, in successive order, whatever the (continuous) proportion
might be.

11.~These magnitudes are said to be corresponding (magnitudes):
the leading to the leading (of two ratios), and
the following to the following.

12.~An alternate ratio is a taking of the (ratio of the) leading  (magnitude) to the leading (of two equal ratios), and (setting it equal to) the (ratio of the)
following   (magnitude) to the following.$^{\P\P}$

13.~An inverse  ratio is a taking of the (ratio of the) following (magnitude) as the leading and the leading (magnitude) as the
following.$^{\ast\ast}$

14.~A composition of a  ratio is a taking of the (ratio of the) leading plus the following (magnitudes),
as one, to the  following (magnitude) by itself.$^{\$\$}$

 15.~A separation of a  ratio is a taking of the (ratio of the) excess by which the
 leading (magnitude) exceeds the following to the following (magnitude)
 by itself.$^{\flat\flat}$
 
16.~A conversion  of a ratio is a taking of the (ratio of the) leading (magnitude) 
 to the excess by which the leading (magnitude) exceeds the following.$^{\dag\dag\dag}$
 
17.~There being several magnitudes,
   and  other (magnitudes)
 of equal  number to them, (which are)  also in
 the same ratio taken two by two,
  a  ratio via equality (i.e.,  ex aequali) occurs when  as the first is to the last  in the first (set of) magnitudes, so  the first (is) to the last   in the second (set of) magnitudes. Or alternately, (it is) a taking of the (ratio of the) outer (magnitudes) by the removal of the 
 inner (magnitudes).$^{\ddag\ddag\ddag}$
 
18.~There being three magnitudes,
  and  other (magnitudes)
 of equal  number to them,  a  perturbed proportion occurs  when  as 
 the  leading is to the following
 in the first (set of) magnitudes, so the  leading (is) to the following  in the second (set of) magnitudes, and as   the following (is) to
 some other ({\rm i.e.}, the remaining magnitude) in the first (set of) magnitudes, so some other (is)  to the leading  in the second (set of) magnitudes.$^{\S\S\S}$}
\end{Parallel}
{\footnotesize \noindent$^\dag$ In other words, $\alpha$ is said to be a part of $\beta$ if $\beta = m\,\alpha$.\\[0.5ex]
$^\ddag$ In modern notation, the ratio of two 
magnitudes, $\alpha$ and $\beta$, is denoted $\alpha:\beta$.\\[0.5ex]
$^\S$ In other words, $\alpha$ has a ratio with respect to $\beta$ if $m\,\alpha >\beta$
and $n\,\beta>\alpha$, for some $m$ and $n$.\\[0.5ex]
$^\P$ In other words, $\alpha:\beta::\gamma:\delta$ if and only if $m\,\alpha>n\,\beta$ whenever $m\,\gamma>n\,\delta$, and $m\,\alpha=n\,\beta$
whenever $m\,\gamma=n\,\delta$, and $m\,\alpha< n\,\beta$ whenever $m\,\gamma<
n\,\delta$, for all $m$ and $n$. This definition  is the kernel of Eudoxus' theory of proportion, and is valid even if $\alpha$,  $\beta$, {\rm etc.}, are irrational.\\[0.5ex]
$^\ast$ Thus if $\alpha$ and $\beta$ have the same ratio
as $\gamma$ and $\delta$ then they are proportional. In modern notation,
$\alpha:\beta::\gamma:\delta$.\\[0.5ex]
$^{\$}$ In modern notation, a proportion in three terms---$\alpha$, $\beta$, and $\gamma$---is written: $\alpha:\beta::\beta:\gamma$.\\
$^{\flat}$ Literally, ``double''.\\[0.5ex]
$^{\dag\dag}$ In other words, if $\alpha:\beta::\beta:\gamma$ then $\alpha:\gamma::\alpha^{\,2}:\beta^{\,2}$.\\[0.5ex]
$^{\ddag\ddag}$ Literally, ``triple''.\\[0.5ex]
$^{\S\S}$ In other words, if $\alpha:\beta::\beta:\gamma::\gamma:\delta$ then $\alpha:\delta::\alpha^{\,3}:\beta^{\,3}$.\\[0.5ex]
$^{\P\P}$ In other words, if $\alpha:\beta::\gamma:\delta$ then the alternate ratio
corresponds to  $\alpha:\gamma::\beta:\delta$.\\[0.5ex]
$^{\ast\ast}$ In other words, if  $\alpha:\beta$ then the inverse ratio
corresponds to $\beta:\alpha$.\\[0.5ex]
$^{\$\$}$ In other words, if
 $\alpha:\beta$ then the composed ratio
 corresponds to  $\alpha+\beta:\beta$.\\[0.5ex]
 $^{\flat\flat}$ In other words, if
 $\alpha:\beta$ then the separated ratio
 corresponds to  $\alpha-\beta:\beta$.\\[0.5ex]
 $^{\dag\dag\dag}$ In other words, if
 $\alpha:\beta$ then the converted ratio
 corresponds to  $\alpha:\alpha-\beta$.\\[0.5ex]
 $^{\ddag\ddag\ddag}$ In other words, if $\alpha, \beta, \gamma$ are the first set of
 magnitudes, and $\delta, \epsilon, \zeta$ the second set, and $\alpha:
 \beta:\gamma::\delta:\epsilon:\zeta$,  then
 the ratio via equality (i.e, ex aequali) corresponds to
 $\alpha:\gamma::\delta:\zeta$.\\[0.5ex]
 $^{\S\S\S}$ In other words, if $\alpha, \beta, \gamma$ are the first set of
 magnitudes, and $\delta, \epsilon, \zeta$ the second set, and $\alpha:\beta::\delta:\epsilon$ as well as $\beta:\gamma::\zeta:\delta$, then
 the proportion is said to be perturbed.}

%%%%%%
% Prop 5.1
%%%%%%
\pdfbookmark[1]{Proposition 5.1}{pdf5.1}
\begin{Parallel}{}{} 
\ParallelLText{
\begin{center}
{\large \ggn{1}.}
\end{center}\vspace*{-7pt}

{\greekfont Ἐὰν ᾖ ὁποσαοῦν μεγέθη ὁποσωνοῦν μεγεθῶν ἴσων τὸ πλῆθος
ἕκαστον ἑκάστου ἰσάκις πολλαπλάσιον, ὁσαπλάσιόν ἐστιν ἓν
τῶν μεγεθῶν ἑνός, τοσαυταπλάσια ἔσται καὶ τὰ πάντα τῶν πάντων.}

{\greekfont Ἔστω ὁποσαοῦν μεγέθη τὰ ΑΒ, ΓΔ ὁποσωνοῦν μεγεθῶν τῶν
Ε, Ζ ἴσων τὸ πλῆθος ἕκαστον ἑκάστου ἰσάκις πολλαπλάσιον; λέγω, ὅτι
ὁσαπλάσιόν ἐστι τὸ ΑΒ τοῦ Ε, τοσαυταπλάσια ἔσται καὶ τὰ ΑΒ, ΓΔ
τῶν Ε, Ζ.}\\~\\~\\

\epsfysize=0.6in
\centerline{\epsffile{Book05/fig01g.eps}}

{\greekfont Ἐπεὶ γὰρ ἰσάκις ἐστὶ πολλαπλάσιον τὸ ΑΒ τοῦ Ε καὶ τὸ ΓΔ τοῦ
Ζ, ὅσα ἄρα ἐστὶν ἐν τῷ ΑΒ μεγέθη ἴσα τῷ Ε, τοσαῦτα καὶ
ἐν τῷ ΓΔ ἴσα τῷ Ζ. διῃρήσθω τὸ μὲν ΑΒ εἰς τὰ τῷ Ε μεγέθη ἴσα τὰ ΑΗ, ΗΒ, τὸ δὲ ΓΔ εἰς τὰ τῷ Ζ ἴσα τὰ ΓΘ, ΘΔ;
ἔσται δὴ ἴσον τὸ πλῆθος τῶν ΑΗ, ΗΒ τῷ πλήθει
τῶν ΓΘ, ΘΔ. καὶ ἐπεὶ ἴσον ἐστὶ τὸ μὲν ΑΗ τῷ Ε, τὸ δὲ 
ΓΘ τῷ Ζ, ἴσον ἄρα τὸ ΑΗ τῷ Ε, καὶ τὰ ΑΗ, ΓΘ τοῖς Ε, Ζ. διὰ
τὰ α ὐτὰ δὴ ἴσον ἐστὶ τὸ ΗΒ τῷ Ε, καὶ τὰ ΗΒ, ΘΔ τοῖς
Ε, Ζ; ὅσα ἄρα ἐστὶν ἐν τῷ ΑΒ ἴσα τῷ Ε, τοσαῦτα καὶ ἐν τοῖς
ΑΒ, ΓΔ ἴσα τοῖς Ε, Ζ; ὁσαπλάσιον ἄρα ἐστὶ τὸ ΑΒ τοῦ Ε,
τοσαυταπλάσια ἔσται καὶ τὰ ΑΒ, ΓΔ τῶν Ε, Ζ.}

{\greekfont Ἐὰν ἄρα ᾖ ὁποσαοῦν μεγέθη ὁποσωνοῦν μεγεθῶν ἴσων τὸ πλῆθος
ἕκαστον ἑκάστου ἰσάκις πολλαπλάσιον, ὁσαπλάσιόν ἐστιν ἓν
τῶν μεγεθῶν ἑνός, τοσαυταπλάσια ἔσται καὶ τὰ πάντα τῶν πάντων;
ὅπερ ἔδει δεῖχαι.}}

\ParallelRText{
\begin{center}
{\large Proposition 1}$^\dag$
\end{center}

If there are any number of magnitudes whatsoever (which are)  equal multiples, respectively, of
some (other) magnitudes,  of equal number (to them), (then) as many times as one of the (first) magnitudes is (divisible) by one (of the second), so many times
will all (of the first magnitudes) also (be divisible) by all (of the second).

Let there be any number of magnitudes whatsoever, $AB$, $CD$, (which are) equal multiples,
respectively, of some (other) magnitudes, $E$, $F$,  of equal  number
(to them). I say that as many times as $AB$ is (divisible) by $E$, so many times will
$AB$, $CD$ also be (divisible) by $E$, $F$.

\epsfysize=0.6in
\centerline{\epsffile{Book05/fig01e.eps}}

For since $AB$, $CD$ are equal multiples of $E$, $F$, thus as many
magnitudes as (there) are in  $AB$ equal to $E$, so many (are there) also in $CD$ equal to $F$. Let $AB$ be divided into magnitudes $AG$, $GB$, equal to $E$,
and $CD$ into (magnitudes) $CH$, $HD$, equal to $F$. So, the number of (divisions)
$AG$, $GB$ will be equal to the number of (divisions) $CH$, $HD$. And since
$AG$ is equal to $E$, and $CH$ to $F$,   $AG$ (is)  thus equal to $E$, and $AG$, $CH$ to
$E$, $F$. So, for the same (reasons), $GB$ is equal to $E$, and $GB$, $HD$ to $E$, $F$.
Thus, as many (magnitudes) as (there) are in $AB$ equal to $E$, so many (are there)
also in $AB$, $CD$ equal to $E$, $F$. Thus, as many times as $AB$ is (divisible) by $E$, so many
times will $AB$, $CD$ also be (divisible)  by $E$, $F$.

Thus, if there are any number of magnitudes whatsoever (which are) equal multiples, respectively, of
some (other) magnitudes,  of  equal  number (to them), (then) as many times as one of the (first) magnitudes is (divisible) by one (of the second), so many times
will all (of the first magnitudes) also (be divisible) by all (of the second). (Which is)
the very thing it was required to show.}
\end{Parallel}
{\footnotesize \noindent$^\dag$ In modern notation, this proposition
reads $m\,\alpha + m\,\beta +\cdots= m\,(\alpha+\beta+\cdots)$.}

%%%%%%
% Prop 5.2
%%%%%%
\pdfbookmark[1]{Proposition 5.2}{pdf5.2}
\begin{Parallel}{}{} 
\ParallelLText{
\begin{center}
{\large \ggn{2}.}
\end{center}\vspace*{-7pt}

{\greekfont Ἐὰν πρῶτον δευτέρου ἰσάκις ᾖ πολλαπλάσιον καὶ τρίτον τετάρτου, ᾖ
δὲ καὶ πέμπτον δευτέρου ἰσάκις πολλαπλάσιον καὶ ἕκτον τετάρτου, καὶ
συντεθὲν πρῶτον καὶ πέμπτον δευτέρου ἰσάκις ἔσται πολλαπλάσιον
καὶ τρίτον καὶ ἕκτον τετάρτου.}

{\greekfont Πρῶτον γὰρ τὸ ΑΒ δευτέρου τοῦ Γ ἰσάκις ἔστω πολλαπλάσιον καὶ
τρίτον τὸ ΔΕ τετάρτου τοῦ Ζ, ἔστω δὲ καὶ πέμπτον τὸ ΒΗ δευτέρου τοῦ
Γ ἰσάκις πολλαπλάσιον καὶ ἕκτον τὸ ΕΘ τετάρτου τοῦ Ζ; λέγω,
ὅτι καὶ συντεθὲν πρῶτον καὶ πέμπτον τὸ ΑΗ δευτέρου τοῦ 
Γ ἰσάκις ἔσται πολλαπλάσιον καὶ τρίτον καὶ ἕκτον
τὸ ΔΘ τετάρτου τοῦ Ζ.}

{\greekfont Ἐπεὶ γὰρ ἰσάκις ἐστὶ πολλαπλάσιον τὸ ΑΒ τοῦ Γ καὶ τὸ 
ΔΕ τοῦ Ζ, ὅσα ἄρα ἐστὶν ἐν τῷ ΑΒ ἴσα τῷ Γ, τοσαῦτα καὶ ἐν
τῷ ΔΕ ἴσα τῷ Ζ. διὰ τὰ α ὐτὰ δὴ καὶ ὅσα ἐστὶν ἐν τῷ ΒΗ
ἴσα τῷ Γ, τοσαῦτα καὶ ἐν τῷ ΕΘ ἴσα τῷ Ζ; ὅσα ἄρα ἐστὶν
ἐν ὅλῳ τῷ ΑΗ ἴσα τῷ Γ, τοσαῦτα καὶ ἐν ὅλῳ τῷ ΔΘ ἴσα τῷ
Ζ; ὁσαπλάσιον ἄρα ἐστὶ τὸ ΑΗ τοῦ Γ, τοσαυταπλάσιον ἔσται καὶ τὸ
ΔΘ τοῦ Ζ. καὶ συντεθὲν ἄρα πρῶτον καὶ πέμπτον τὸ ΑΗ δευτέρου
τοῦ Γ ἰσάκις ἔσται πολλαπλάσιον καὶ τρίτον καὶ ἕκτον τὸ
ΔΘ τετάρτου τοῦ Ζ.}\\~\\~\\~\\~\\~\\~\\~\\~\\

\epsfysize=1.5in
\centerline{\epsffile{Book05/fig02g.eps}}

{\greekfont Ἐὰν ἄρα πρῶτον δευτέρου ἰσάκις ᾖ πολλαπλάσιον καὶ τρίτον τετάρτου, ᾖ
δὲ καὶ πέμπτον δευτέρου ἰσάκις πολλαπλάσιον καὶ ἕκτον τετάρτου, καὶ
συντεθὲν πρῶτον καὶ πέμπτον δευτέρου ἰσάκις ἔσται πολλαπλάσιον
καὶ τρίτον καὶ ἕκτον τετάρτου; ὅπερ ἔδει δεῖχαι.}}

\ParallelRText{
\begin{center}
{\large Proposition 2}$^{\dag}$
\end{center}

If a first (magnitude) and a third are equal multiples of a second
and a fourth (respectively), and a fifth (magnitude) and a sixth (are) also equal multiples of the
second and fourth (respectively), (then) the first (magnitude) and the fifth, being added together, and the third and the sixth,  (being added together), will also be equal multiples of the
second (magnitude) and the fourth (respectively).

For let a first (magnitude) $AB$ and a third $DE$ be equal multiples of a
second $C$ and a fourth $F$ (respectively). And let a fifth (magnitude) $BG$
and a sixth $EH$ also be (other) equal multiples of the second $C$ and the fourth $F$
(respectively). I say that the first (magnitude) and the fifth, being added together, (to give) $AG$,
and the third (magnitude) and the sixth, (being added together, to give) $DH$, will also be
equal multiples of the second (magnitude) $C$ and the fourth $F$ (respectively).

For since $AB$ and $DE$ are equal multiples of $C$ and $F$ (respectively), 
thus as many (magnitudes) as (there) are in $AB$ equal to $C$, so many (are there) also
in $DE$ equal to $F$.  And so, for the same (reasons),   as many (magnitudes) as (there) are in $BG$ equal to $C$, so many (are there) also in $EH$ equal to $F$.
Thus, as many (magnitudes) as (there) are in the whole of $AG$ equal to $C$,
so many (are there) also in the whole of $DH$ equal to $F$.
Thus, as many times as $AG$ is (divisible) by $C$, so many times will $DH$ also
be divisible by $F$. Thus, the first (magnitude) and the fifth, being
added together, (to give) $AG$, and the third (magnitude) and the sixth, (being
added together, to give) $DH$, will also be equal multiples of the second (magnitude)
$C$ and the fourth $F$ (respectively).

\epsfysize=1.5in
\centerline{\epsffile{Book05/fig02e.eps}}

Thus, if a first (magnitude) and a third are equal multiples of a second
and a fourth (respectively), and a fifth (magnitude) and a sixth (are) also equal multiples of the
second and fourth (respectively), (then) the  first (magnitude) and the fifth, being
added together, and the third  and sixth, (being added together), will also be equal multiples of the
second (magnitude) and the fourth (respectively). (Which is) the very thing it was required to
show.}
\end{Parallel}
{\footnotesize \noindent$^\dag$ In modern notation, this propostion reads
$m\,\alpha+n\,\alpha = (m+n)\,\alpha$.}

%%%%%%
% Prop 5.3
%%%%%%
\pdfbookmark[1]{Proposition 5.3}{pdf5.3}
\begin{Parallel}{}{} 
\ParallelLText{
\begin{center}
{\large \ggn{3}.}
\end{center}\vspace*{-7pt}

{\greekfont Ἐὰν πρῶτον δευτέρου ἰσάκις ᾖ πολλαπλάσιον καὶ τρίτον τετάρτου,
ληφθῇ δὲ ἰσάκις πολλαπλάσια τοῦ τε πρώτου καὶ τρίτου, καὶ δι 
ἴσου τῶν  ληφθέντων ἑκάτερον ἑκατέρου ἰσάκις ἔσται πολλαπλάσιον
τὸ μὲν τοῦ δευτέρου τὸ δὲ τοῦ τετάρτου.}

{\greekfont Πρῶτον γὰρ τὸ Α δευτέρου τοῦ Β ἰσάκις ἔστω πολλαπλάσιον καὶ
τρίτον τὸ Γ τετάρτου τοῦ Δ, καὶ εἰλήφθω τῶν Α, Γ ἰσάκις πολλαπλάσια
τὰ ΕΖ, ΗΘ; λέγω, ὅτι ἰσάκις ἐστὶ πολλαπλάσιον τὸ ΕΖ τοῦ Β καὶ
τὸ ΗΘ τοῦ Δ.}

{\greekfont Ἐπεὶ γὰρ ἰσάκις ἐστὶ πολλαπλάσιον τὸ ΕΖ τοῦ Α καὶ τὸ ΗΘ τοῦ
Γ, ὅσα ἄρα ἐστὶν ἐν τῷ ΕΖ ἴσα τῷ Α, τοσαῦτα καὶ ἐν τῷ
ΗΘ ἴσα τῷ Γ. διῃρήσθω τὸ μὲν ΕΖ εἰς τὰ τῷ Α μεγέθη ἴσα
τὰ ΕΚ, ΚΖ, τὸ δὲ ΗΘ εἰς τὰ τῷ Γ ἴσα τὰ ΗΛ, ΛΘ; ἔσται δὴ ἴσον
τὸ πλῆθος τῶν ΕΚ, ΚΖ τῷ πλήθει τῶν ΗΛ, ΛΘ. καὶ ἐπεὶ
ἰσάκις ἐστὶ πολλαπλάσιον τὸ Α τοῦ Β καὶ τὸ Γ τοῦ Δ,  ἴσον δὲ τὸ
μὲν ΕΚ τῷ Α, τὸ δὲ ΗΛ τῷ Γ, ἰσάκις ἄρα ἐστὶ πολλαπλάσιον
τὸ ΕΚ τοῦ Β καὶ τὸ ΗΛ τοῦ Δ. διὰ τὰ α ὐτὰ δὴ ἰσάκις ἐστὶ
πολλαπλάσιον τὸ ΚΖ τοῦ Β καὶ τὸ ΛΘ τοῦ Δ. ἐπεὶ οῦ̓ν
πρῶτον τὸ ΕΚ δευτέρου τοῦ Β ἴσάκις ἐστὶ πολλαπλάσιον καὶ
τρίτον τὸ ΗΛ τετάρτου τοῦ Δ, ἔστι δὲ καὶ πέμπτον τὸ ΚΖ δευτέρου τοῦ
Β ἰσάκις πολλαπλάσιον καὶ ἕκτον τὸ ΛΘ τετάρτου τοῦ Δ,
καὶ συντεθὲν  ἄρα πρῶτον καὶ πέμπτον τὸ ΕΖ δευτέρου τοῦ Β
ἰσάκις ἐστὶ πολλαπλάσιον καὶ τρίτον καὶ ἕκτον τὸ ΗΘ
τετάρτου τοῦ Δ.}\\~\\~\\~\\~\\~\\

\epsfysize=2in
\centerline{\epsffile{Book05/fig03g.eps}}

{\greekfont Ἐὰν ἄρα πρῶτον δευτέρου ἰσάκις ᾖ πολλαπλάσιον καὶ τρίτον τετάρτου,
ληφθῇ δὲ τοῦ πρώτου καὶ τρίτου ἰσάκις πολλαπλάσια, καὶ δι 
ἴσου τῶν  ληφθέντων ἑκάτερον ἑκατέρου ἰσάκις ἔσται πολλαπλάσιον
τὸ μὲν τοῦ δευτέρου τὸ δὲ τοῦ τετάρτου; ὅπερ ἔδει δεῖχαι.}}

\ParallelRText{
\begin{center}
{\large Proposition 3}$^\dag$
\end{center}

If a first (magnitude) and a third are equal multiples
of a second and a fourth (respectively), and equal multiples are taken of the
first and the third, (then), via equality,  the 
(magnitudes) taken will also be equal multiples of the second (magnitude) and the fourth, respectively.

For let a first (magnitude) $A$ and a third $C$ be equal multiples of a
second $B$ and a fourth $D$ (respectively), and let the equal multiples
$EF$ and $GH$ be taken of $A$ and $C$ (respectively). I say that
$EF$ and $GH$ are equal multiples of $B$ and $D$ (respectively).

For since $EF$ and $GH$ are equal multiples of $A$ and $C$ (respectively), thus as many (magnitudes) as (there) are in $EF$ equal to $A$, so many (are there) also in
$GH$ equal to $C$. Let $EF$ be divided into magnitudes $EK$, $KF$ equal
to $A$, and $GH$ into (magnitudes) $GL$, $LH$ equal to $C$. So, the number of (magnitudes)
$EK$, $KF$ will be equal to the number of (magnitudes) $GL$, $LH$.
And since $A$ and $C$ are equal multiples of $B$ and $D$ (respectively), and $EK$ (is)
equal to $A$, and $GL$ to $C$,  $EK$ and $GL$ are thus equal multiples of $B$ and $D$
(respectively). So, for the same (reasons), $KF$ and $LH$ are equal multiples
of $B$ and $D$ (respectively). Therefore, since the first (magnitude) $EK$ and the
third $GL$ are equal multiples of the second $B$ and the fourth $D$ (respectively), and the fifth (magnitude) $KF$ and
the sixth $LH$ are  also  equal multiples of the second $B$ and the fourth $D$ (respectively),  then  the first (magnitude) and fifth, being added together,
(to give) $EF$, and the third (magnitude) and sixth, (being added together, to give) $GH$, are thus also equal multiples of the second (magnitude) $B$ and the fourth $D$ (respectively) [Prop.~5.2].

\epsfysize=2in
\centerline{\epsffile{Book05/fig03e.eps}}

Thus, if a first (magnitude) and a third are equal multiples
of a second and a fourth (respectively), and equal multiples are taken of the
first and the third, (then), via equality,  the 
(magnitudes) taken will also be equal multiples of the second (magnitude) and the fourth, respectively. (Which is) the very thing it was required to show.}
\end{Parallel}
{\footnotesize \noindent$^\dag$ In modern notation, this proposition reads
$m\,(n\,\alpha) = (m\,n)\,\alpha$.}


%%%%%%
% Prop 5.4
%%%%%%
\pdfbookmark[1]{Proposition 5.4}{pdf5.4}
\begin{Parallel}{}{} 
\ParallelLText{
\begin{center}
{\large \ggn{4}.}
\end{center}\vspace*{-7pt}

{\greekfont Ἐὰν πρῶτον πρὸς δεύτερον τὸν α ὐτὸν ἔθῃ λόγον καὶ τρίτον
πρὸς τέταρτον, καὶ τὰ ἰσάκις πολλαπλάσια τοῦ τε πρώτου καὶ
τρίτου πρὸς τὰ ἰσάκις πολλαπλάσια τοῦ δευτέρου καὶ τετάρτου καθ
ὁποιονοῦν πολλαπλασιασμὸν τὸν α ὐτὸν ἕχει λόγον ληφθέντα
κατάλληλα.}

{\greekfont Πρῶτον γὰρ τὸ Α πρὸς δεύτερον τὸ Β τὸν α ὐτὸν ἐθέτω λόγον
καὶ τρίτον τὸ Γ πρὸς τέταρτον τὸ Δ, καὶ εἰλήφθω τῶν μὲν Α, Γ
ἰσάκις πολλαπλάσια τὰ Ε, Ζ, τῶν δὲ Β, Δ ἄλλα, ἃ ἔτυθεν, ἰσάκις
πολλαπλάσια τὰ Η, Θ; λέγω, ὅτι ἐστὶν ὡς τὸ Ε πρὸς τὸ Η, οὕτως 
τὸ Ζ πρὸς τὸ Θ.}

{\greekfont Εἰλήφθω γὰρ τῶν μὲν Ε, Ζ ἰσάκις πολλαπλάσια τὰ Κ, Λ, τῶν δὲ Η,
Θ ἄλλα, ἃ ἔτυθεν, ἰσάκις πολλαπλάσια τὰ Μ, Ν.}

{\greekfont μ̀βοχ{[}Καὶ] ἐπεὶ ἰσάκις ἐστὶ πολλαπλάσιον τὸ μὲν Ε τοῦ Α, τὸ δὲ Ζ τοῦ
Γ, καὶ εἴληπται τῶν Ε, Ζ ἴσάκις πολλαπλάσια τὰ Κ, Λ, ἴσάκις
ἄρα ἐστὶ πολλαπλάσιον τὸ Κ τοῦ Α καὶ τὸ Λ τοῦ Γ. διὰ τὰ α ὐτὰ δὴ
ἰσάκις ἐστὶ πολλαπλάσιον τὸ Μ τοῦ Β καὶ τὸ Ν τοῦ Δ. καὶ ἐπεί
ἐστιν ὡς τὸ Α πρὸς τὸ Β, οὕτως τὸ Γ πρὸς τὸ Δ, καὶ εἴληπται
τῶν μὲν Α, Γ ἰσάκις πολλαπλάσια τὰ Κ, Λ, τῶν δὲ Β, Δ ἄλλα, ἃ
ἔτυθεν, ἰσάκις πολλαπλάσια τὰ Μ, Ν, εἰ ἄρα ὑπερέθει τὸ Κ
τοῦ Μ, ὑπερέθει καὶ τὸ Λ τοῦ Ν, καὶ εἰ ἴσον, ἴσον, καὶ
εἰ ἔλαττον, ἔλαττον. καί ἐστι τὰ μὲν Κ, Λ τῶν Ε, Ζ ἰσάκις 
πολλαπλάσια, τὰ δὲ Μ, Ν τῶν Η, Θ ἄλλα, ἃ ἔτυθεν, ἰσάκις
πολλαπλάσια;
 ἔστιν ἄρα ὡς τὸ Ε πρὸς τὸ Η, οὕτως τὸ Ζ
πρὸς τὸ Θ.}\\~\\~\\~\\

\epsfysize=3in
\centerline{\epsffile{Book05/fig04g.eps}}

{\greekfont Ἐὰν ἄρα πρῶτον πρὸς δεύτερον τὸν α ὐτὸν ἔθῃ λόγον καὶ τρίτον
πρὸς τέταρτον, καὶ τὰ ἰσάκις πολλαπλάσια τοῦ τε πρώτου καὶ
τρίτου πρὸς τὰ ἰσάκις πολλαπλάσια τοῦ δευτέρου καὶ τετάρτου τὸν α ὐτὸν ἕχει λόγον καθ ὁποιονοῦν πολλαπλασιασμὸν  ληφθέντα
κατάλληλα; ὅπερ ἔδει δεῖχαι.}}

\ParallelRText{
\begin{center}
{\large Proposition 4}$^\dag$
\end{center}

If  a first (magnitude) has the same ratio to a second 
that a third  (has) to a fourth, (then) equal multiples of the first (magnitude) and the third will also have the same ratio to equal multiples of the second and the fourth,  being taken in corresponding order, according to any kind of
multiplication whatsoever.

For let a first (magnitude) $A$ have the same ratio to a second $B$ that
a third $C$  (has) to a fourth $D$. And let equal multiples $E$ and $F$ 
be taken of $A$ and $C$ (respectively), and  other random equal multiples
$G$ and $H$  of $B$ and $D$ (respectively). I say that as $E$ (is)
to $G$, so $F$ (is) to $H$.

For let equal multiples $K$ and $L$ be taken of $E$ and $F$ (respectively),
and other random equal multiples $M$ and $N$ of $G$ and $H$ (respectively).

\mbox{[}And] since $E$ and $F$ are equal multiples of $A$ and $C$ (respectively), and
the equal multiples $K$ and $L$ have been taken of $E$ and $F$ (respectively),
 $K$ and $L$ are thus equal multiples of $A$ and $C$ (respectively) [Prop.~5.3]. So, for the same (reasons), $M$ and $N$
 are equal multiples of $B$ and $D$ (respectively). And since as $A$ is to
 $B$, so $C$ (is) to $D$, and the equal multiples $K$ and $L$ have been taken of $A$ and
 $C$ (respectively), and the other random equal multiples  $M$ and $N$ of $B$ and $D$ (respectively), (then) if $K$ exceeds $M$ (then) $L$ also
 exceeds $N$, and if ($K$ is) equal (to $M$ then $L$ is also) equal (to $N$), and if ($K$ is) less (than $M$ then $L$ is also) less (than $N$) [Def.~5.5]. And $K$ and $L$ are equal multiples of 
 $E$ and $F$ (respectively), and $M$ and $N$ other random equal multiples of $G$ and $H$ (respectively). Thus, as $E$ (is) to $G$, so $F$ (is) to $H$ 
 [Def.~5.5].

\epsfysize=3in
\centerline{\epsffile{Book05/fig04e.eps}}

Thus, if a first (magnitude) has the same ratio to a second 
that a third  (has) to a fourth, (then) equal multiples of the first (magnitude) and the third will also have the same ratio to equal multiples of the second and the fourth,  being taken in corresponding order, according to any kind of
multiplication whatsoever. (Which is) the very thing
it was required to show.}
\end{Parallel}
{\footnotesize \noindent$^\dag$ In modern notation, this
proposition reads that if $\alpha:\beta :: \gamma:\delta$ then
$m\,\alpha:n\,\beta::m\,\gamma:n\,\delta$, for all $m$ and $n$.}

%%%%%%
% Prop 5.5
%%%%%%
\pdfbookmark[1]{Proposition 5.5}{pdf5.5}
\begin{Parallel}{}{} 
\ParallelLText{
\begin{center}
{\large \ggn{5}.}
\end{center}\vspace*{-7pt}

{\greekfont Ἐὰν μέγεθος μεγέθους ἰσάκις ᾖ πολλαπλάσιον, ὅπερ ἀφαιρεθὲν
ἀφαιρεθέντος, καὶ τὸ λοιπὸν τοῦ λοιποῦ ἰσάκις ἔσται πολλαπλάσιον,
ὁσαπλάσιόν ἐστι τὸ ὅλον τοῦ ὅλου.}\\~\\

\epsfysize=0.7in
\centerline{\epsffile{Book05/fig05g.eps}}

{\greekfont Μέγεθος γὰρ τὸ ΑΒ μεγέθους τοῦ ΓΔ ἰσάκις ἔστω πολλαπλάσιον,
ὅπερ ἀφαιρεθὲν τὸ ΑΕ ἀφαιρεθέντος τοῦ ΓΖ; λέγω, ὅτι καὶ λοιπὸν
τὸ ΕΒ λοιποῦ τοῦ ΖΔ ἰσάκις ἔσται πολλαπλάσιον, ὁσαπλάσιόν 
ἐστιν ὅλον τὸ ΑΒ ὅλου τοῦ ΓΔ.}

{\greekfont Ὁσαπλάσιον γάρ ἐστι τὸ ΑΕ τοῦ ΓΖ, τοσαυταπλάσιον γεγονέτω καὶ
τὸ ΕΒ τοῦ ΓΗ.}

{\greekfont Καὶ ἐπεὶ ἰσάκις ἐστὶ πολλαπλάσιον τὸ ΑΕ τοῦ ΓΖ καὶ τὸ ΕΒ τοῦ
ΗΓ, ἰσάκις ἄρα ἐστὶ πολλαπλάσιον τὸ ΑΕ τοῦ ΓΖ καὶ τὸ ΑΒ τοῦ
ΗΖ. κεῖται δὲ ἰσάκις πολλαπλάσιον τὸ ΑΕ τοῦ ΓΖ καὶ τὸ ΑΒ τοῦ
ΓΔ. ἰσάκις ἄρα ἐστὶ πολλαπλάσιον τὸ ΑΒ ἑκατέρου τῶν ΗΖ, ΓΔ;
ἴσον ἄρα τὸ ΗΖ τῷ ΓΔ. κοινὸν ἀφῃρήσθω τὸ ΓΖ; λοιπὸν
ἄρα τὸ ΗΓ λοιπῷ τῷ ΖΔ ἴσον ἐστίν. καὶ ἐπεὶ ἰσάκις
ἐστὶ πολλαπλάσιον τὸ ΑΕ τοῦ ΓΖ καὶ τὸ ΕΒ τοῦ ΗΓ, ἴσον
δὲ τὸ ΗΓ τῷ  ΔΖ, 
ἰσάκις ἄρα  ἐστὶ πολλαπλάσιον τὸ ΑΕ τοῦ ΓΖ καὶ τὸ ΕΒ τοῦ ΖΔ.
ἰσάκις δὲ ὑπόκειται πολλαπλάσιον τὸ ΑΕ τοῦ ΓΖ καὶ τὸ ΑΒ τοῦ
ΓΔ; ἰσάκις ἄρα 
ἐστὶ πολλαπλάσιον τὸ ΕΒ τοῦ ΖΔ καὶ τὸ ΑΒ τοῦ ΓΔ. καὶ λοιπὸν 
ἄρα τὸ ΕΒ λοιποῦ τοῦ ΖΔ ἰσάκις ἔσται πολλαπλάσιον, ὁσαπλάσιόν
ἐστιν ὅλον τὸ ΑΒ ὅλου τοῦ ΓΔ.}

{\greekfont Ἐὰν ἄρα μέγεθος μεγέθους ἰσάκις ᾖ πολλαπλάσιον, ὅπερ ἀφαιρεθὲν
ἀφαιρεθέντος, καὶ τὸ λοιπὸν τοῦ λοιποῦ ἰσάκις ἔσται πολλαπλάσιον,
ὁσαπλάσιόν ἐστι καὶ τὸ ὅλον τοῦ ὅλου; ὅπερ ἔδει δεῖχαι.}}

\ParallelRText{
\begin{center}
{\large Proposition 5}$^\dag$
\end{center}

If a magnitude is the same multiple of a magnitude
that a (part) taken away (is) of a (part) taken away (respectively), (then) the remainder will
also be the same multiple of the remainder as that  which the whole (is) of the whole (respectively).

\epsfysize=0.7in
\centerline{\epsffile{Book05/fig05e.eps}}

For let the magnitude $AB$ be the same multiple of the magnitude $CD$ that
the (part)  taken away $AE$ (is) of the (part)  taken away $CF$ (respectively).
I say that the remainder $EB$ will also be the same multiple
of the remainder $FD$ as that which the whole $AB$ (is) of the whole
$CD$ (respectively).

For as many times as $AE$ is (divisible) by $CF$, so many times let $EB$
also be made (divisible) by $CG$.

And since $AE$ and $EB$ are equal multiples of $CF$ and $GC$ (respectively),
$AE$ and $AB$ are thus equal multiples of $CF$ and $GF$ (respectively) 
 [Prop.~5.1]. And  $AE$ and $AB$ are assumed (to be) equal multiples of $CF$ and $CD$ (respectively). Thus,
 $AB$ is an equal multiple of each of $GF$ and $CD$. Thus, $GF$ (is) equal to
 $CD$. Let $CF$ be subtracted from both. Thus, the remainder $GC$ is
 equal to the remainder $FD$. And since $AE$ and $EB$ are equal multiples
 of $CF$ and $GC$ (respectively), and $GC$ (is) equal to $DF$, $AE$  and $EB$ are thus
 equal multiples of $CF$ and $FD$ (respectively). And $AE$ and $AB$ are assumed
 (to be) equal multiples of $CF$ and $CD$ (respectively). Thus, $EB$ and $AB$ are
 equal multiples of $FD$ and $CD$ (respectively). Thus, the remainder
 $EB$ will also be the same multiple of the remainder $FD$
 as that which the whole $AB$ (is) of the whole $CD$ (respectively).
 
 Thus, if a magnitude is the same multiple of a magnitude
that a (part) taken away (is) of a (part) taken away (respectively), (then) the remainder will
also be the same multiple of the remainder as that  which the whole (is) of the whole (respectively).
  (Which is) the very thing it was required to show.}
\end{Parallel}
{\footnotesize \noindent$^\dag$ In modern notation, this
proposition reads $m\,\alpha-m\,\beta=m\,(\alpha-\beta)$.}

%%%%%%
% Prop 5.6
%%%%%%
\pdfbookmark[1]{Proposition 5.6}{pdf5.6}
\begin{Parallel}{}{} 
\ParallelLText{
\begin{center}
{\large \ggn{6}.}
\end{center}\vspace*{-7pt}

{\greekfont Ἐὰν δύο μεγέθη δύο μεγεθῶν ἰσάκις ᾖ πολλαπλάσια, καὶ ἀφαιρεθέντα
τινὰ τῶν α ὐτῶν ἰσάκις ᾖ πολλαπλάσια, καὶ τὰ λοιπὰ τοῖς α ὐτοῖς
ἤτοι ἴσα ἐστὶν ἢ ἰσάκις α ὐτῶν πολλαπλάσια.}

{\greekfont Δύο γὰρ μεγέθη τὰ ΑΒ, ΓΔ δύο μεγεθῶν τῶν Ε, Ζ ἰσάκις ἔστω
πολλαπλάσια, καὶ ἀφαιρεθέντα τὰ ΑΗ, ΓΘ τῶν α ὐτῶν τῶν Ε, Ζ ἰσάκις
ἔστω πολλαπλάσια;  λέγω, ὅτι καὶ λοιπὰ τὰ ΗΒ, ΘΔ τοῖς Ε, Ζ
ἤτοι ἴσα ἐστὶν ἢ ἰσάκις α ὐτῶν πολλαπλάσια.}\\~\\~\\

\epsfysize=1.8in
\centerline{\epsffile{Book05/fig06g.eps}}

{\greekfont Ἔστω γὰρ πρότερον τὸ ΗΒ τῷ Ε ἴσον; λέγω, ὅτι καὶ τὸ ΘΔ τῷ Ζ
ἴσον ἐστίν.}

{\greekfont Κείσθω γὰρ τῷ Ζ ἴσον τὸ ΓΚ. ἐπεὶ ἰσάκις ἐστὶ πολλαπλάσιον
τὸ ΑΗ τοῦ Ε καὶ τὸ ΓΘ τοῦ Ζ,  ἴσον δὲ τὸ μὲν ΗΒ τῷ Ε,
τὸ δὲ ΚΓ τῷ Ζ, ἰσάκις ἄρα ἐστὶ πολλαπλάσιον τὸ ΑΒ τοῦ Ε καὶ
τὸ ΚΘ τοῦ Ζ. ἰσάκις δὲ ὑπόκειται πολλαπλάσιον τὸ ΑΒ τοῦ Ε καὶ
τὸ ΓΔ τοῦ Ζ; ἴσάκις ἄρα ἐστὶ πολλαπλάσιον τὸ ΚΘ τοῦ Ζ καὶ τὸ ΓΔ τοῦ Ζ. ἐπεὶ οὖν ἑκάτερον τῶν ΚΘ, ΓΔ
τοῦ Ζ ἰσάκις ἐστὶ πολλαπλάσιον, ἴσον ἄρα ἐστὶ τὸ ΚΘ τῷ ΓΔ.
κοινὸν ἀφῃρήσθω τὸ ΓΘ; λοιπὸν ἄρα τὸ ΚΓ λοιπῷ τῷ ΘΔ
ἴσον ἐστίν. ἀλλὰ τὸ Ζ τῷ ΚΓ ἐστιν ἴσον; καὶ τὸ ΘΔ ἄρα
τῷ Ζ ἴσον ἐστίν. ὥστε εἰ τὸ ΗΒ τῷ Ε ἴσον ἐστίν, καὶ τὸ ΘΔ
ἴσον ἔσται τῷ Ζ.}

{\greekfont Ὁμοίως δὴ δείχομεν, ὅτι, κᾂν πολλαπλάσιον ᾖ τὸ ΗΒ τοῦ Ε,
τοσαυταπλάσιον ἔσται καὶ τὸ ΘΔ τοῦ Ζ.}

{\greekfont Ἐὰν ἄρα δύο μεγέθη δύο μεγεθῶν ἰσάκις ᾖ πολλαπλάσια, καὶ ἀφαιρεθέντα
τινὰ τῶν α ὐτῶν ἰσάκις ᾖ πολλαπλάσια, καὶ τὰ λοιπὰ τοῖς α ὐτοῖς
ἤτοι ἴσα ἐστὶν ἢ ἰσάκις α ὐτῶν πολλαπλάσια; ὅπερ ἔδει
δεῖχαι.}}

\ParallelRText{
\begin{center}
{\large Proposition 6}$^\dag$
\end{center}

If two magnitudes are equal multiples of two (other) magnitudes, and some (parts) taken away (from the former magnitudes) are equal multiples of the latter (magnitudes, respectively),
(then) the remainders  are also either equal to the latter (magnitudes), or (are) equal multiples of them (respectively).

For let two magnitudes $AB$ and $CD$ be equal multiples of two magnitudes
$E$ and $F$ (respectively). And let the (parts)  taken away (from the former) $AG$ and $CH$ be equal multiples of $E$ and $F$ (respectively). I say that the remainders $GB$ and
$HD$ are also either equal to $E$ and $F$ (respectively), or (are) equal multiples of them.

\epsfysize=1.8in
\centerline{\epsffile{Book05/fig06e.eps}}

For let $GB$ be, first of all, equal to $E$. I say that $HD$ is also equal to $F$.

For let  $CK$ be made equal to $F$. Since $AG$ and $CH$ are equal multiples of $E$ and $F$ (respectively), and $GB$ (is) equal to $E$, and $KC$ to $F$,
$AB$ and $KH$ are thus equal multiples of $E$ and $F$ (respectively)
[Prop.~5.2].
And $AB$ and $CD$ are assumed (to be) equal multiples of $E$ and $F$ (respectively). Thus, $KH$ and $CD$ are equal multiples of $F$ and $F$ (respectively). Therefore, $KH$ and $CD$ are each equal multiples of $F$. Thus, $KH$
is equal to $CD$. Let $CH$  be taken away from both. Thus, the remainder
$KC$ is equal to the remainder $HD$. But, $F$ is equal to $KC$. Thus, $HD$ is also equal to $F$. Hence, if $GB$ is equal to $E$, (then) $HD$ will also be equal to $F$.

So, similarly, we can show that even if $GB$ is a multiple of $E$, (then) $HD$
will also be the same multiple of $F$.

Thus, if two magnitudes are equal multiples of two (other) magnitudes, and some (parts) taken away (from the former magnitudes) are equal multiples of the latter (magnitudes, respectively),
(then) the remainders  are also either equal to the latter (magnitudes), or (are) equal multiples of them (respectively). (Which is) the very thing it was required to show.}
\end{Parallel}
{\footnotesize \noindent$^\dag$ In modern notation, this
proposition reads $m\,\alpha-n\,\alpha = (m-n)\,\alpha$.}

%%%%%%
% Prop 5.7
%%%%%%
\pdfbookmark[1]{Proposition 5.7}{pdf5.7}
\begin{Parallel}{}{} 
\ParallelLText{
\begin{center}
{\large \ggn{7}.}
\end{center}\vspace*{-7pt}

{\greekfont Τὰ ἴσα πρὸς τὸ α ὐτὸ τὸν α ὐτὸν ἔχει λόγον καὶ τὸ α ὐτὸ πρὸς
τὰ ἴσα.}

{\greekfont Ἔστω ἴσα μεγέθη τὰ Α, Β, ἄλλο δέ τι, ὃ ἔτυθεν,
μέγεθος τὸ Γ; λέγω, ὅτι ἑκάτερον τῶν Α, Β πρὸς τὸ Γ τὸν
α ὐτὸν ἔχει λόγον, καὶ τὸ Γ πρὸς ἑκάτερον τῶν Α, Β.}\\

\epsfysize=0.75in
\centerline{\epsffile{Book05/fig07g.eps}}

{\greekfont Εἰλήφθω γὰρ τῶν μὲν Α, Β ἰσάκις πολλαπλάσια τὰ Δ, Ε, τοῦ δὲ Γ
ἄλλο, ὃ ἔτυθεν, πολλαπλάσιον τὸ Ζ.}

{\greekfont Ἐπεὶ οὖν ἰσάκις ἐστὶ πολλαπλάσιον τὸ Δ τοῦ Α καὶ τὸ Ε τοῦ Β,
ἴσον δὲ τὸ Α τῷ Β, ἴσον ἄρα καὶ τὸ Δ τῷ Ε. ἄλλο δέ, ὅ ἔτυθεν,
τὸ Ζ. εἰ ἄρα ὑπερέθει τὸ Δ τοῦ Ζ, ὑπερέθει καὶ τὸ Ε τοῦ Ζ, καὶ
εἰ ἴσον, ἴσον, καὶ εἰ ἔλαττον, ἔλαττον. καί ἐστι τὰ μὲν Δ, Ε
τῶν Α, Β ἰσάκις πολλαπλάσια, τὸ δὲ Ζ τοῦ Γ ἄλλο, ὃ ἔτυθεν,
πολλαπλάσιον; ἔστιν ἄρα ὡς τὸ Α πρὸς τὸ Γ, οὕτως τὸ Β πρὸς τὸ
Γ.}

{\greekfont Λέγω [δή], ὅτι καὶ τὸ Ε πρὸς ἑκάτερον τῶν Α, Β τὸν α ὐτὸν
ἔχει λόγον.}

{\greekfont Τῶν γὰρ α ὐτῶν κατασκευασθέντων ὁμοίως δείχομεν, ὅτι
ἴσον ἐστὶ τὸ Δ τῷ Ε; ἄλλο δέ τι τὸ Ζ;
εἰ ἄρα ὑπερέθει τὸ Ζ τοῦ Δ, ὑπερέθει καὶ τοῦ Ε, καὶ εἰ ἴσον,
ἴσον, καὶ εἰ ἔλαττον, ἔλαττον. καί ἐστι τὸ μὲν Ζ τοῦ Γ
πολλαπλάσιον, τὰ δὲ Δ, Ε τῶν Α, Β ἄλλα, ἃ ἔτυθεν, ἰσάκις
πολλαπλάσια; ἔστιν ἄρα ὡς τὸ Γ πρὸς τὸ Α, οὕτως τὸ Γ πρὸς
τὸ Β.}

{\greekfont Τὰ ἴσα ἄρα πρὸς τὸ α ὐτὸ τὸν α ὐτὸν ἔχει λόγον καὶ τὸ α ὐτὸ πρὸς
τὰ ἴσα.}\\~\\~\\~\\

\begin{center}
{\large {\greekfont Πόρισμα}.}
\end{center}\vspace*{-7pt}

{\greekfont Ἐκ δὴ τούτου φανερόν, ὅτι ἐὰν μεγέθη τινὰ ἀνάλογον ᾖ, καὶ
ἀνάπαλιν ἀνάλογον ἔσται. ὅπερ ἔδει δεῖχαι.}}

\ParallelRText{
\begin{center}
{\large Proposition 7}
\end{center}

Equal (magnitudes) have the same ratio to the same 
(magnitude), and the latter (magnitude has the same ratio) to the equal
(magnitudes).

Let $A$ and $B$ be equal magnitudes, and $C$ some other random magnitude.
I say that $A$ and $B$ each have the same ratio to $C$, and (that) $C$ (has the same
ratio) to each of $A$ and $B$.

\epsfysize=0.75in
\centerline{\epsffile{Book05/fig07e.eps}}

For let the equal multiples $D$ and $E$ be taken of $A$ and $B$ (respectively), and the other
random multiple $F$ of $C$.

Therefore, since $D$ and $E$ are equal multiples of $A$ and $B$ (respectively), and
$A$ (is) equal to $B$, $D$ (is) thus also equal to $E$. And $F$ (is) different, at random.
Thus, if $D$ exceeds $F$ (then) $E$ also exceeds $F$, and if ($D$ is) equal (to $F$ then
$E$ is also) equal (to $F$), and if ($D$ is) less (than $F$ then
$E$ is also) less (than $F$). And $D$ and $E$ are equal multiples of $A$ and $B$ (respectively), and $F$ another random multiple of $C$. Thus, as $A$ (is) to
$C$, so $B$ (is) to $C$ [Def.~5.5].

\mbox{[}So] I say that $C$$^\dag$ also has the same ratio to each of $A$ and $B$.

For, similarly,  we can show, by the same construction, that $D$ is equal to $E$.
And $F$ (has) some other (value). Thus, if $F$ exceeds $D$ then it also exceeds $E$,
and if ($F$ is) equal (to $D$  then it is also) equal (to $E$), and if ($F$ is) less (than $D$ then it is also) less (than $E$). And $F$ is a multiple of $C$, and $D$ and $E$ other random equal
multiples of $A$ and $B$. Thus, as $C$ (is) to $A$, so $C$ (is)
to $B$  [Def.~5.5]. 

Thus, equal (magnitudes) have the same ratio to the same 
(magnitude), and the latter (magnitude has the same ratio) to the equal
(magnitudes).\\

\begin{center}
\large{Corollary}$^\ddag$
\end{center}\vspace*{-7pt}

So (it is) clear, from this, that if some magnitudes are proportional,
(then) they will also be proportional inversely. (Which is) the very thing it was
required to show.} 
\end{Parallel}
{\footnotesize \noindent$^\dag$ The Greek text has ``$E$'', which is obviously a mistake.\\[0.5ex]
$^\ddag$ In modern notation, this corollary reads that if $\alpha:\beta::\gamma:\delta$ then
$\beta:\alpha::\delta:\gamma$.}

%%%%%%
% Prop 5.8
%%%%%%
\pdfbookmark[1]{Proposition 5.8}{pdf5.8}
\begin{Parallel}{}{} 
\ParallelLText{
\begin{center}
{\large \ggn{8}.}
\end{center}\vspace*{-7pt}

{\greekfont Τῶν ἀνίσων μεγεθῶν τὸ μεῖζον πρὸς τὸ α ὐτὸ μείζονα λόγον ἔχει ἤπερ τὸ ἔλαττον. καὶ τὸ α ὐτὸ πρὸς τὸ ἔλαττον μείζονα
λόγον ἔχει ἤπερ πρὸς τὸ μεῖζον.}

{\greekfont Ἔστω ἄνισα μεγέθη τὰ ΑΒ, Γ, καὶ ἔστω μεῖζον τὸ ΑΒ, ἄλλο δέ,
ὃ ἔτυθεν, τὸ Δ; λέγω, ὅτι τὸ ΑΒ πρὸς τὸ Δ μείζονα λόγον
ἔχει ἤπερ τὸ Γ πρὸς τὸ Δ, καὶ τὸ Δ πρὸς τὸ Γ μείζονα λόγον
ἔχει ἤπερ πρὸς τὸ ΑΒ.}\\

\epsfysize=1.8in
\centerline{\epsffile{Book05/fig08g.eps}}

{\greekfont Ἐπεὶ γὰρ μεῖζόν ἐστι τὸ ΑΒ τοῦ Γ, κείσθω τῷ Γ ἴσον τὸ ΒΕ;
τὸ δὴ ἔλασσον τῶν ΑΕ, ΕΒ πολλαπλασιαζόμενον ἔσται ποτὲ τοῦ
Δ μεῖζον. ἔστω πρότερον τὸ ΑΕ ἔλαττον τοῦ ΕΒ, καὶ πεπολλαπλασιάσθω
τὸ ΑΕ, καὶ ἔστω α ὐτοῦ πολλαπλάσιον τὸ ΖΗ μεῖζον ὂν τοῦ
Δ, καὶ ὁσαπλάσιόν ἐστι τὸ ΖΗ τοῦ ΑΕ, τοσαυταπλάσιον γεγονέτω
καὶ τὸ μὲν ΗΘ τοῦ ΕΒ τὸ δὲ Κ τοῦ Γ; καὶ εἰλήφθω τοῦ Δ 
διπλάσιον μὲν τὸ Λ, τριπλάσιον δὲ τὸ Μ, καὶ ἑχῆς ἑνὶ πλεῖον,
ἕως ἂν τὸ λαμβανόμενον πολλαπλάσιον μὲν γένηται τοῦ Δ, 
πρώτως δὲ μεῖζον τοῦ Κ. εἰλήφθω, καὶ ἔστω τὸ Ν τετραπλάσιον
μὲν τοῦ Δ, πρώτως δὲ μεῖζον τοῦ Κ.}

{\greekfont Ἐπεὶ οὖν τὸ Κ τοῦ Ν  πρώτως ἐστὶν ἔλαττον, τὸ Κ ἄρα τοῦ 
Μ οὔκ ἐστιν ἔλαττον. καὶ ἐπεὶ ἰσάκις ἐστὶ πολλαπλάσιον
τὸ ΖΗ τοῦ ΑΕ καὶ τὸ ΗΘ τοῦ ΕΒ, ἰσάκις ἄρα ἐστὶ πολλαπλάσιον
τὸ ΖΗ τοῦ ΑΕ καὶ τὸ ΖΘ τοῦ ΑΒ. ἰσάκις δέ ἐστι πολλαπλάσιον
τὸ ΖΗ τοῦ ΑΕ καὶ τὸ Κ τοῦ Γ; ἰσάκις ἄρα ἐστὶ πολλαπλάσιον
τὸ ΖΘ τοῦ ΑΒ καὶ τὸ Κ τοῦ Γ. τὰ ΖΘ, Κ ἄρα τῶν ΑΒ, Γ ἰσάκις
ἐστὶ πολλαπλάσια. πάλιν, ἐπεὶ ἰσάκις ἐστὶ πολλαπλάσιον τὸ ΗΘ 
τοῦ ΕΒ καὶ τὸ Κ τοῦ Γ, ἴσον δὲ τὸ ΕΒ τῷ Γ, ἴσον ἄρα
καὶ τὸ ΗΘ τῷ Κ. τὸ δὲ Κ τοῦ Μ οὔκ ἐστιν ἔλαττον; οὐδ
ἄρα τὸ ΗΘ τοῦ Μ ἔλαττόν ἐστιν. μεῖζον δὲ τὸ ΖΗ τοῦ Δ; ὅλον
ἄρα τὸ ΖΘ συναμφοτέρων τῶν Δ, Μ μεῖζόν ἐστιν. ἀλλὰ συναμφότερα
τὰ Δ, Μ τῷ Ν ἐστιν ἴσα, ἐπειδήπερ τὸ Μ τοῦ Δ τριπλάσιόν
ἐστιν, συναμφότερα δὲ τὰ Μ, Δ τοῦ Δ ἐστι τετραπλάσια, ἔστι
δὲ καὶ τὸ Ν τοῦ Δ τετραπλάσιον; συναμφότερα ἄρα τὰ Μ, Δ τῷ Ν ἴσα
ἐστίν. ἀλλὰ τὸ ΖΘ τῶν Μ, Δ μεῖζόν ἐστιν; τὸ ΖΘ ἄρα τοῦ Ν
ὑπερέθει; τὸ δὲ Κ τοῦ Ν οὐθ ὑπερέθει. καί ἐστι τὰ μὲν ΖΘ, Κ τῶν ΑΒ, Γ ἰσάκις πολλαπλάσια, τὸ δὲ Ν τοῦ Δ ἄλλο, ὃ
ἔτυθεν, πολλαπλάσιον; τὸ ΑΒ ἄρα πρὸς τὸ Δ μείζονα λόγον
ἔχει ἤπερ τὸ Γ πρὸς τὸ Δ.}

{\greekfont Λέγω δή, ὅτι καὶ τὸ Δ πρὸς τὸ Γ μείζονα λόγον ἔχει ἤπερ
τὸ Δ πρὸς τὸ ΑΒ.}

{\greekfont Τῶν γὰρ α ὐτῶν κατασκευασθέντων ὁμοίως δείχομεν, ὅτι
τὸ μὲν Ν τοῦ Κ ὑπερέθει, τὸ δὲ Ν τοῦ ΖΘ οὐθ ὑπερέθει.
καί ἐστι τὸ μὲν Ν τοῦ Δ πολλαπλάσιον, τὰ δὲ ΖΘ, Κ τῶν ΑΒ, Γ
ἄλλα, ἃ ἔτυθεν, ἰσάκις πολλαπλάσια; τὸ Δ ἄρα πρὸς τὸ Γ
μείζονα λόγον ἔχει ἤπερ τὸ Δ πρὸς τὸ ΑΒ.}

{\greekfont Ἀλλὰ δὴ τὸ ΑΕ τοῦ ΕΒ μεῖζον ἔστω. τὸ δὴ ἔλαττον τὸ ΕΒ
πολλαπλασιαζόμενον ἔσται ποτὲ τοῦ Δ μεῖζον. πεπολλαπλασιάσθω, καὶ
ἔστω τὸ ΗΘ πολλαπλάσιον μὲν τοῦ ΕΒ, μεῖζον δὲ τοῦ Δ; καὶ
ὁσαπλάσιόν ἐστι τὸ ΗΘ τοῦ ΕΒ, τοσαυταπλάσιον γεγονέτω
καὶ τὸ μὲν ΖΗ τοῦ ΑΕ, τὸ δὲ Κ τοῦ Γ. ὁμοίως δὴ δείχομεν,
ὅτι τὰ ΖΘ, Κ τῶν ΑΒ, Γ ἰσάκις ἐστὶ πολλαπλάσια; καὶ εἰλήφθω
ὁμοίως τὸ Ν πολλαπλάσιον μὲν τοῦ Δ, πρώτως δὲ μεῖζον
τοῦ ΖΗ; ὥστε πάλιν τὸ ΖΗ τοῦ Μ οὔκ ἐστιν ἔλασσον. μεῖζον
δὲ τὸ ΗΘ τοῦ Δ; ὅλον ἄρα τὸ ΖΘ τῶν Δ, Μ, τουτέστι τοῦ Ν,
ὑπερέθει. τὸ δὲ Κ τοῦ Ν οὐθ ὑπερέθει, ἐπειδήπερ καὶ τὸ
ΖΗ μεῖζον ὂν τοῦ ΗΘ, τουτέστι τοῦ Κ, τοῦ Ν οὐθ ὑπερέθει.
καὶ ὡσαύτως κατακολουθοῦντες τοῖς ἐπάνω περαίνομεν
τὴν ἀπόδειχιν.}

{\greekfont Τῶν ἄρα ἀνίσων μεγεθῶν τὸ μεῖζον πρὸς τὸ α ὐτὸ μείζονα λόγον ἔχει ἤπερ τὸ ἔλαττον; καὶ τὸ α ὐτὸ πρὸς τὸ ἔλαττον μείζονα
λόγον ἔχει ἤπερ πρὸς τὸ μεῖζον; ὅπερ ἔδει δεῖχαι.}}

\ParallelRText{
\begin{center}
{\large Proposition 8}
\end{center}

For unequal magnitudes, the greater (magnitude) has a greater
ratio  than the lesser to the same (magnitude). And the latter (magnitude)
has a greater ratio to the lesser (magnitude) than to the greater.

Let $AB$ and $C$ be unequal magnitudes, and let $AB$ be the greater (of the two),
and $D$ another random magnitude. I say that $AB$ has a greater ratio to $D$
than  $C$ (has) to $D$, and (that) $D$ has a greater ratio to $C$ than (it has) to $AB$.

\epsfysize=1.8in
\centerline{\epsffile{Book05/fig08e.eps}}

For since $AB$ is greater than $C$, let $BE$ be made equal to $C$.
So, the lesser of $AE$ and $EB$, being multiplied, will sometimes be greater than $D$ [Def.~5.4]. First of all, let $AE$ be less than $EB$, and let $AE$ be multiplied, 
and let $FG$ be a multiple of it which (is) greater than $D$.
And as many times as $FG$ is (divisible) by $AE$, so many times
let $GH$ also  become (divisible) by $EB$, and $K$ by $C$. And
let the double multiple $L$ of $D$ be taken, and the triple
multiple $M$, and several more, (each increasing) in order by one,
until the (multiple) taken  becomes the first multiple of $D$ (which is) greater than $K$. Let it be taken, and let it also be the quadruple multiple
$N$ of $D$---the first (multiple) greater than $K$.

Therefore, since $K$ is less than $N$ first, $K$ is thus not less than $M$.
And since $FG$ and $GH$ are equal multiples of $AE$ and $EB$ (respectively),
$FG$ and $FH$ are thus equal multiples of $AE$ and $AB$ (respectively) [Prop.~5.1].
And $FG$ and $K$ are equal multiples of $AE$ and $C$ (respectively).
Thus, $FH$ and $K$ are equal multiples of $AB$ and $C$ (respectively).
Thus, $FH$, $K$ are equal multiples of $AB$, $C$. Again, since $GH$ and 
$K$ are equal
multiples of $EB$ and $C$, and $EB$ (is) equal to $C$, $GH$ (is) thus also equal to
$K$. And $K$ is not less than $M$. Thus, $GH$ not less than $M$ either. And $FG$
(is) greater than $D$. Thus, the whole of $FH$ is greater than $D$
and $M$ (added) together. But, $D$ and $M$ (added) together is equal to $N$, 
inasmuch as $M$ is three times $D$, and $M$ and $D$ (added) together is four
times $D$, and $N$ is also four times $D$. Thus, $M$ and $D$ (added) together is
equal to $N$. But, $FH$ is greater than $M$ and $D$. Thus, $FH$ exceeds $N$. And $K$ does not exceed $N$. 
And $FH$, $K$  are equal multiples of $AB$, $C$, and $N$ another random multiple
of $D$. Thus, $AB$ has a greater ratio to $D$  than $C$ (has) to $D$ [Def.~5.7].

So, I say that $D$ also has a greater ratio to $C$ than $D$ (has) to $AB$.

For, similarly, by the same construction, we can show that 
$N$ exceeds $K$, and $N$ does not exceed $FH$. And $N$ is a multiple of $D$,
and $FH$, $K$ other  random  equal multiples of $AB$, $C$ (respectively). Thus, $D$ has a greater
ratio to  $C$ than $D$ (has) to $AB$ [Def.~5.5].

And so let $AE$ be greater than $EB$. So, the lesser, $EB$, being multiplied, will
sometimes be greater than $D$. Let it be multiplied, and let $GH$ be
a multiple of $EB$ (which is) greater than $D$. And as many times as $GH$ is
(divisible) by $EB$, so many times let $FG$ also  become (divisible) by $AE$,
and $K$ by $C$. So, similarly (to the above), we can show that $FH$ and $K$ are equal multiples of
$AB$ and $C$ (respectively). And, similarly (to the above), let the multiple $N$ of $D$,  (which is)
the first (multiple) greater than $FG$, be taken. So,  $FG$ is again not less than $M$.
And $GH$ (is) greater than $D$. Thus, the whole of $FH$ exceeds $D$ and $M$, that is to say $N$.  And $K$ does not exceed $N$, inasmuch as $FG$, which (is) greater than
$GH$---that is to say, $K$---also does not exceed $N$. And, following the above (arguments), we  (can) complete the proof in the same manner.

Thus, for unequal magnitudes, the greater (magnitude) has a greater
ratio  than the lesser to the same (magnitude). And the latter (magnitude)
has a greater ratio to the lesser (magnitude) than to the greater. (Which is)
the very thing it was required to show.}
\end{Parallel}

%%%%%%
% Prop 5.9
%%%%%%
\pdfbookmark[1]{Proposition 5.9}{pdf5.9}
\begin{Parallel}{}{} 
\ParallelLText{
\begin{center}
{\large \ggn{9}.}
\end{center}\vspace*{-7pt}

{\greekfont Τὰ πρὸς τὸ α ὐτὸ τὸν α ὐτὸν ἔθοντα λὸγον ἴσα ἀλλήλοις ἐστίν;
καὶ πρὸς ἃ τὸ α ὐτὸ τὸν α ὐτὸν ἕθει λόγον, ἐκεῖνα ἴσα ἐστίν.}

\epsfysize=0.5in
\centerline{\epsffile{Book05/fig09g.eps}}

{\greekfont Ἐθέτω γὰρ ἑκάτερον τῶν Α, Β πρὸς τὸ Γ τὸν α ὐτὸν λόγον;
λέγω, ὅτι ἴσον ἐστὶ τὸ Α τῷ Β.}

{\greekfont Εἰ γὰρ μή, οὐκ ἂν ἑκάτερον τῶν Α, Β πρὸς τὸ Γ τὸν α ὐτὸν
εἶθε λόγον; ἔχει δέ; ἴσον ἄρα ἐστὶ τὸ Α τῷ Β.}

{\greekfont Ἐθέτω δὴ πάλιν τὸ Γ πρὸς ἑκάτερον τῶν Α, Β τὸν α ὐτὸν
λόγον; λέγω, ὅτι ἴσον ἐστὶ τὸ Α τῷ Β.}

{\greekfont Εἰ γὰρ μή, οὐκ ἂν τὸ Γ πρὸς ἑκάτερον τῶν Α, Β τὸν α ὐτὸν εἶθε λόγον; ἔχει δέ; ἴσον ἄρα ἐστὶ τὸ Α τῷ Β.}

{\greekfont Τὰ ἄρα πρὸς τὸ α ὐτὸ τὸν α ὐτὸν ἔθοντα λόγον ἴσα ἀλλήλοις
ἐστίν; καὶ πρὸς ἃ τὸ α ὐτὸ τὸν α ὐτὸν ἔχει λόγον, ἐκεῖνα
ἴσα ἐστίν; ὅπερ ἔδει δεῖχαι.}}

\ParallelRText{
\begin{center}
{\large Proposition 9}
\end{center}

(Magnitudes) having the same ratio to the same (magnitude) are equal to one another. And those (magnitudes) to which the same (magnitude)
has the same ratio are equal.

\epsfysize=0.5in
\centerline{\epsffile{Book05/fig09e.eps}}

For let $A$ and $B$ each have the same ratio to $C$. I say that $A$ is equal to $B$.

For if not, $A$ and $B$ would not each have the same ratio to $C$ [Prop.~5.8].
But they do. Thus, $A$ is equal to $B$.

So, again, let $C$ have the same ratio to each of $A$ and $B$. I say that $A$ is equal to $B$.

For if not, $C$ would not have the same ratio to each of $A$ and $B$ [Prop.~5.8]. But it does. Thus, $A$ is equal to $B$.

Thus, (magnitudes) having the same ratio to the same (magnitude) are equal to one another. And those (magnitudes) to which the same (magnitude)
has the same ratio are equal. (Which is) the very thing it was required to show.}
\end{Parallel}

%%%%%%
% Prop 5.10
%%%%%%
\pdfbookmark[1]{Proposition 5.10}{pdf5.10}
\begin{Parallel}{}{} 
\ParallelLText{
\begin{center}
{\large \ggn{10}.}
\end{center}\vspace*{-7pt}

{\greekfont Τῶν πρὸς τὸ α ὐτὸ λόγον ἐθόντων τὸ μείζονα λόγον ἔθον ἐκεῖνο μεῖζόν ἐστιν; πρὸς ὃ δὲ τὸ α ὐτὸ μείζονα λόγον ἔχει, ἐκεῖνο
ἔλαττόν ἐστιν.}\\

\epsfysize=0.5in
\centerline{\epsffile{Book05/fig10g.eps}}

{\greekfont Ἐθέτω γὰρ τὸ Α πρὸς τὸ Γ μείζονα λόγον ἤπερ τὸ Β πρὸς τὸ Γ;
λέγω, ὅτι μεῖζόν ἐστι τὸ Α  τοῦ Β.}

{\greekfont Εἰ γὰρ μή, ἤτοι ἴσον ἐστὶ τὸ Α τῷ Β ἢ ἔλασσον. ἴσον
μὲν οὖν οὔκ ἐστὶ τὸ Α τῷ Β; ἑκάτερον γὰρ  ἂν τῶν Α, Β
πρὸς τὸ Γ τὸν α ὐτὸν εἶθε λόγον. οὐκ ἔχει δέ; οὐκ ἄρα ἴσον
ἐστὶ τὸ Α τῷ Β. οὐδὲ μὴν ἔλασσόν ἐστι τὸ Α τοῦ Β; τὸ Α γὰρ
ἂν πρὸς τὸ Γ  ἐλάσσονα λόγον εἶθεν ἤπερ τὸ Β πρὸς τὸ Γ. οὐκ
ἔχει δέ; οὐκ ἄρα ἔλασσόν ἐστι τὸ Α τοῦ Β. ἐδείθθη δὲ οὐδὲ ἴσον; μεῖζον ἄρα ἐστὶ τὸ Α τοῦ Β.}

{\greekfont Ἐθέτω δὴ πάλιν τὸ Γ πρὸς τὸ Β μείζονα λόγον ἤπερ τὸ Γ πρὸς
τὸ Α; λέγω, ὅτι ἔλασσόν ἐστι τὸ Β τοῦ Α.}

{\greekfont Εἰ γὰρ μή, ἤτοι ἴσον ἐστὶν ἢ μεῖζον. ἴσον μὲν οὖν οὔκ
ἐστι τὸ Β τῷ Α; τὸ Γ γὰρ ἂν πρὸς ἑκάτερον τῶν Α, Β τὸν α ὐτὸν
εἶθε λόγον. οὐκ ἔχει δέ; οὐκ ἄρα ἴσον ἐστὶ τὸ Α τῷ Β. οὐδὲ
μὴν μεῖζόν ἐστι τὸ Β τοῦ Α; τὸ Γ γὰρ ἂν πρὸς τὸ Β ἐλάσσονα
λόγον εἶθεν ἤπερ πρὸς τὸ Α. οὐκ ἔχει δέ; οὐκ ἄρα
μεῖζόν ἐστι τὸ Β τοῦ Α. ἐδείθθη δέ, ὅτι οὐδὲ ἴσον; ἔλαττον
ἄρα ἐστὶ τὸ Β τοῦ Α.}

{\greekfont Τῶν ἄρα  πρὸς τὸ α ὐτὸ λόγον ἐθόντων τὸ μείζονα λόγον ἔθον μεῖζόν ἐστιν; καὶ πρὸς ὃ  τὸ α ὐτὸ μείζονα λόγον ἔχει, ἐκεῖνο
ἔλαττόν ἐστιν; ὅπερ ἔδει δεῖχαι.}}

\ParallelRText{
\begin{center}
{\large Proposition 10}
\end{center}

For (magnitudes) having a ratio to the same (magnitude), that (magnitude which) has the
greater ratio is (the) greater. And that (magnitude) to which the latter (magnitude) has a greater ratio is (the)
lesser.

\epsfysize=0.5in
\centerline{\epsffile{Book05/fig10e.eps}}

For let $A$ have a greater ratio to $C$ than $B$ (has) to $C$. I say that $A$ is greater than $B$.

For if not, $A$ is surely either equal to or less than $B$. In fact, $A$ is not
equal to $B$. For (then) $A$ and $B$ would each have the same ratio to $C$ 
 [Prop.~5.7]. But they do not. Thus, $A$
 is not equal to $B$. Neither, indeed, is $A$ less than $B$. For (then)
 $A$ would have a lesser ratio to $C$ than $B$ (has) to $C$ [Prop.~5.8]. But it does not. Thus, $A$ is not less
 than $B$. And it was shown not (to be) equal either. Thus, $A$ is greater than $B$.
 
 So, again, let $C$ have a greater ratio to $B$ than $C$ (has) to $A$. I say that
 $B$ is less than $A$.
 
 For if not, (it is) surely either equal or greater. In fact, $B$ is not equal to 
 $A$. For (then) $C$ would have the same ratio to each of $A$ and $B$  [Prop.~5.7]. But it does not. Thus, $A$ is not equal to $B$. 
 Neither, indeed, is $B$ greater than $A$. For (then) $C$ would have a lesser
 ratio to $B$ than (it has) to $A$ [Prop.~5.8].
 But it does not. Thus, $B$ is not greater than $A$. And it was shown that
  (it is) not equal (to $A$) either. Thus, $B$ is less than $A$.
  
 Thus, for (magnitudes) having a ratio to the same (magnitude), that (magnitude which)
 has the greater  ratio is (the) greater. And that (magnitude) to which the latter (magnitude) has a greater ratio is (the)
lesser. (Which is) the very thing it was required to show.}
\end{Parallel}

%%%%%%
% Prop 5.11
%%%%%%
\pdfbookmark[1]{Proposition 5.11}{pdf5.11}
\begin{Parallel}{}{} 
\ParallelLText{
\begin{center}
{\large \ggn{11}.}
\end{center}\vspace*{-7pt}

{\greekfont Οἱ τῷ α ὐτῷ λόγῳ οἱ α ὐτοὶ καὶ ἀλλήλοις εἰσὶν οἱ α ὐτοί.}

{\greekfont Ἔστωσαν γὰρ ὡς μὲν τὸ Α πρὸς τὸ Β, οὕτως τὸ Γ πρὸς τὸ Δ, ὡς
δὲ τὸ Γ πρὸς τὸ Δ, οὕτως τὸ Ε πρὸς τὸ Ζ; λέγω, ὅτι ἐστὶν ὡς
τὸ Α πρὸς τὸ Β, οὕτως τὸ Ε πρὸς τὸ Ζ.}

{\greekfont Εἰλήφθω γὰρ τῶν Α, Γ, Ε ἰσάκις πολλαπλάσια τὰ Η, Θ, Κ, τῶν δὲ Β, Δ, Ζ 
ἄλλα, ἃ ἔτυθεν, ἰσάκις πολλαπλάσια τὰ Λ, Μ, Ν.}

\epsfysize=0.75in
\centerline{\epsffile{Book05/fig11g.eps}}

{\greekfont Καὶ ἐπεί ἐστιν ὡς τὸ Α πρὸς τὸ Β, οὕτως τὸ Γ
πρὸς τὸ Δ, καὶ εἴληπται τῶν μὲν Α, Γ ἰσάκις πολλαπλάσια τὰ Η, Θ,
τῶν δὲ Β, Δ ἄλλα, ἃ ἔτυθεν, ἰσάκις πολλαπλάσια τὰ Λ, Μ, εἰ
ἄρα ὑπερέθει τὸ Η τοῦ Λ, ὑπερέθει καὶ τὸ Θ τοῦ Μ, καὶ εἰ
ἴσον ἐστίν, ἴσον, καὶ εἰ ἐλλείπει, ἐλλείπει. πάλιν, ἐπεί
ἐστιν ὡς τὸ Γ πρὸς τὸ Δ, οὕτως τὸ Ε πρὸς τὸ Ζ, καὶ εἴληπται
τῶν  Γ, Ε ἰσάκις πολλαπλάσια τὰ Θ, Κ, τῶν δὲ Δ, Ζ ἄλλα, ἃ ἔτυθεν,
ἰσάκις πολλαπλάσια τὰ Μ, Ν, εἰ ἄρα ὑπερέθει τὸ Θ τοῦ Μ,
ὑπερέθει καὶ τὸ Κ τοῦ Ν, καὶ εἰ ἴσον, ἴσον, καὶ εἰ
ἔλλατον, ἔλαττον. ἀλλὰ εἰ ὑπερεῖθε τὸ Θ τοῦ Μ, ὑπερεῖθε
καὶ τὸ Η τοῦ Λ, καὶ εἰ ἴσον, ἴσον, καὶ εἰ ἔλαττον, 
ἔλαττον; ὥστε καὶ εἰ ὑπερέθει τὸ Η τοῦ Λ, ὑπερέθει καὶ τὸ
Κ τοῦ Ν, καὶ εἰ ἴσον, ἴσον, καὶ εἰ ἔλαττον, ἔλαττον. καί
ἐστι τὰ  μὲν Η, Κ τῶν Α, Ε ἰσάκις πολλαπλάσια, τὰ δὲ Λ, Ν τῶν
Β, Ζ ἄλλα, ἃ ἔτυθεν, ἰσάκις πολλαπλάσια; ἔστιν ἄρα ὡς τὸ Α πρὸς
τὸ Β, οὕτως τὸ Ε πρὸς τὸ Ζ.}

{\greekfont Οἱ ἄρα τῷ α ὐτῷ λόγῳ οἱ α ὐτοὶ καὶ ἀλλήλοις εἰσὶν οἱ α ὐτοί;
ὅπερ ἔδει δεῖχαι.}}

\ParallelRText{
\begin{center}
{\large Proposition 11}$^\dag$
\end{center}

(Ratios which are) the same with the same ratio  are also the same with one another.

For let it be that as $A$ (is) to $B$, so $C$ (is) to $D$, and as $C$ (is) to $D$, so $E$ (is) to $F$. I say
that as  $A$ is to $B$, so $E$ (is) to $F$.

For let the equal multiples $G$, $H$,  $K$ be taken of $A$, $C$,  $E$ (respectively), and
the other random equal multiples $L$, $M$,  $N$ of $B$, $D$, $F$ (respectively).

\epsfysize=0.75in
\centerline{\epsffile{Book05/fig11e.eps}}

And since as $A$ is to $B$, so $C$ (is) to $D$, and the equal multiples $G$ and $H$
be taken of $A$ and $C$ (respectively), and the other random equal
multiples $L$ and $M$ of $B$ and $D$ (respectively),  thus if $G$ exceeds $L$ (then)
$H$ also exceeds $M$, and if ($G$ is) equal (to $L$ then $H$ is also)
equal (to $M$), and if ($G$ is) less (than $L$ then $H$ is also) less (than $M$)  [Def.~5.5]. Again, since as $C$ is to $D$, so $E$
(is) to $F$, and the equal multiples $H$ and $K$ have been taken of $C$ and $E$
(respectively), and the other random equal multiples $M$ and $N$ of
$D$ and $F$ (respectively), thus if $H$ exceeds $M$ (then)
$K$ also exceeds $N$, and if ($H$ is) equal  (to $M$ then $K$ is also)
equal (to $N$), and if ($H$ is) less (than $M$  then $K$ is also) less (than $N$) [Def.~5.5]. But (we saw that) if $H$ was exceeding $M$ (then) $G$ was also
exceeding $L$, and if ($H$ was) equal (to $M$ then $G$ was also)
equal (to $L$), and if ($H$ was) less (than $M$ then $G$ was also) less (than $L$).
And, hence, if $G$ exceeds $L$ (then)
$K$ also exceeds $N$, and if ($G$ is) equal (to $L$ then $K$ is also)  equal (to $N$), and if ($G$ is) less (than $L$ then $K$ is also) less (than $N$). And $G$ and $K$ are equal multiples of $A$ and $E$ (respectively), and $L$ and $N$
other random equal multiples of $B$ and $F$ (respectively). Thus, as $A$ is to $B$,
so $E$ (is) to $F$ [Def.~5.5].

Thus, (ratios which are) the same with the same ratio  are also the same with one another. (Which is) the very thing it was required to show.}
\end{Parallel}
{\footnotesize \noindent$^\dag$ In modern notation, this proposition
reads that if $\alpha:\beta::\gamma:\delta$ and $\gamma:\delta::\epsilon:\zeta$ then $\alpha:\beta::\epsilon:\zeta$.}

%%%%%%
% Prop 5.12
%%%%%%
\pdfbookmark[1]{Proposition 5.12}{pdf5.12}
\begin{Parallel}{}{} 
\ParallelLText{
\begin{center}
{\large \ggn{12}.}
\end{center}\vspace*{-7pt}

{\greekfont Ἐὰν ᾖ ὁποσαοῦν μεγέθη  ἀνάλογον, ἔσται ὡς ἓν τῶν ἡγουμένων πρὸς ἓν τῶν ἑπομένων, οὕτως ἅπαντα τὰ ἡγούμενα πρὸς ἅπαντα τὰ ἑπόμενα.}\\

\epsfysize=1.3in
\centerline{\epsffile{Book05/fig12g.eps}}

{\greekfont Ἔστωσαν ὁποσαοῦν μεγέθη ἀνάλογον τὰ Α, Β, Γ, Δ, Ε, Ζ, ὡς τὸ
Α πρὸς τὸ Β, οὕτως τὸ Γ πρὸς τὸ Δ, καὶ τὸ Ε πρὸς το Ζ; λέγω, ὅτι
ἐστὶν ὡς τὸ Α πρὸς τὸ Β, οὕτως τὰ Α, Γ, Ε πρὸς τὰ Β, Δ, Ζ.}

{\greekfont Εἰλήφθω γὰρ τῶν μὲν Α, Γ, Ε ἰσάκις πολλαπλάσια τὰ Η, Θ, Κ, τῶν
δὲ Β, Δ, Ζ ἄλλα, ἃ ἔτυθεν, ἰσάκις πολλαπλάσια τὰ Λ, Μ, Ν.}

{\greekfont Καὶ ἐπεί ἐστιν ὡς τὸ Α πρὸς τὸ Β, οὕτως τὸ Γ πρὸς τὸ Δ, καὶ
τὸ Ε πρὸς τὸ Ζ, καὶ εἴληπται τῶν μὲν Α, Γ, Ε ἰσάκις πολλαπλάσια
τὰ Η, Θ, Κ τῶν δὲ Β, Δ, Ζ ἄλλα, ἃ ἔτυθεν, ἰσάκις πολλαπλάσια
τὰ Λ, Μ, Ν, εἰ ἄρα ὑπερέθει τὸ Η τοῦ Λ, ὑπερέθει καὶ τὸ Θ
τοῦ Μ, καὶ τὸ Κ τοῦ Ν, καὶ εἰ ἴσον, ἴσον, καὶ εἰ ἔλαττον, ἔλαττον.
ὥστε καὶ εἰ ὑπερέθει τὸ Η τοῦ Λ, ὑπερέθει καὶ τὰ Η, Θ, Κ τῶν
Λ, Μ, Ν,  καὶ εἰ ἴσον, ἴσα, καὶ εἰ ἔλαττον, ἔλαττονα. καί ἐστι
τὸ μὲν Η καὶ τὰ Η, Θ, Κ τοῦ Α καὶ τῶν Α, Γ, Ε ἰσάκις
πολλαπλάσια, ἐπειδήπερ ἐὰν ᾖ  ὁποσαοῦν μεγέθη ὁποσωνοῦν
μεγεθῶν ἴσων τὸ πλῆθος ἕκαστον ἑκάστου ἰσάκις πολλαπλάσιον,
ὁσαπλάσιόν ἐστιν ἓν τῶν μεγεθῶν ἑνός, τοσαυταπλάσια ἔσται
καὶ τὰ πάντα τῶν πάντων. διὰ τὰ α ὐτὰ δὴ καὶ τὸ Λ καὶ τὰ Λ, Μ, Ν
τοῦ Β καὶ τῶν Β, Δ, Ζ ἰσάκις ἐστὶ πολλαπλάσια;
ἔστιν ἄρα ὡς τὸ Α πρὸς τὸ Β, οὕτως τὰ Α, Γ, Ε πρὸς τὰ Β, Δ, Ζ.}

{\greekfont Ἐὰν ἄρα ᾖ ὁποσαοῦν μεγέθη  ἀνάλογον, ἔσται ὡς ἓν τῶν ἡγουμένων πρὸς ἓν τῶν ἑπομένων, οὕτως ἅπαντα τὰ ἡγούμενα πρὸς ἅπαντα τὰ ἑπόμενα; ὅπερ ἔδει δεῖχαι.}}

\ParallelRText{
\begin{center}
{\large Proposition 12}$^\dag$
\end{center}

If there are any number  of magnitudes whatsoever
(which are) proportional, (then) as one of the leading (magnitudes is) to one of the following, so
will 
all of the leading (magnitudes) be to all of the following.

\epsfysize=1.3in
\centerline{\epsffile{Book05/fig12e.eps}}

Let there be any number of magnitudes whatsoever, 
$A$, $B$, $C$, $D$, $E$, $F$, (which are) proportional, (so that) as $A$ (is) to $B$, so $C$ (is) to $D$, and 
$E$ to $F$. I say that as $A$ is to $B$, so  $A$, $C$, $E$ (are) to  $B$, $D$, $F$.

For let the equal multiples $G$, $H$,  $K$ be taken of $A$, $C$, $E$
(respectively), and the other random equal multiples $L$, $M$,  $N$ of
$B$, $D$, $F$ (respectively).

And since as $A$ is to $B$, so $C$ (is) to $D$, and $E$ to $F$, and the equal multiples
$G$, $H$, $K$ have been taken of $A$, $C$,  $E$ (respectively), and the other
random equal multiples $L$, $M$,  $N$ of $B$, $D$,  $F$ (respectively), thus
if $G$ exceeds $L$ (then) $H$ also exceeds $M$, and $K$ (exceeds) $N$, and if ($G$ is)
equal (to $L$ then $H$ is also) equal (to $M$, and $K$  to $N$), and if ($G$ is) less (than $L$ then $H$ is also) less (than $M$, and $K$  than $N$)  [Def.~5.5].  And, hence, if $G$ exceeds $L$ (then)
$G$, $H$, $K$ also exceed $L$, $M$, $N$, and if ($G$ is) equal
(to $L$ then $G$, $H$, $K$ are also) equal (to  $L$, $M$, $N$) and if ($G$ is) less (than $L$ then
$G$, $H$, $K$ are also) less (than $L$, $M$, $N$). And $G$ and $G$, $H$, $K$ are equal multiples
of $A$ and $A$, $C$, $E$ (respectively), inasmuch as if there are any number of magnitudes whatsoever (which are)  equal multiples, respectively, of
some  (other) magnitudes,  of equal  number (to them), (then) as many times as one of the (first) magnitudes is (divisible) by one (of the second), so many times
will all (of the first magnitudes) also (be divisible) by all (of the second)
 [Prop.~5.1]. So, for the same (reasons), $L$
 and $L$, $M$, $N$ are also equal multiples of $B$ and $B$, $D$, $F$ (respectively). Thus,
 as $A$ is to $B$, so $A$, $C$, $E$ (are) to $B$, $D$, $F$ (respectively).
 
Thus, if there are any number  of magnitudes whatsoever
(which are) proportional, (then) as one of the leading (magnitudes is) to one of the following, so
will 
all of the leading (magnitudes) be to all of the following. (Which is) the
very thing it was required to show.}
\end{Parallel}
{\footnotesize \noindent$^\dag$ In modern notation, this proposition
reads that if $\alpha:\alpha'::\beta:\beta'::\gamma:\gamma'$ {\rm etc.}\ then
$\alpha:\alpha'::(\alpha+\beta+\gamma+\cdots):(\alpha'+\beta'+\gamma'+\cdots)$.}

%%%%%%
% Prop 5.13
%%%%%%
\pdfbookmark[1]{Proposition 5.13}{pdf5.13}
\begin{Parallel}{}{} 
\ParallelLText{
\begin{center}
{\large \ggn{13}.}
\end{center}\vspace*{-7pt}

{\greekfont Ἐὰν πρῶτον πρὸς δεύτερον τὸν α ὐτὸν ἒθῃ  λόγον καὶ τρίτον
πρὸς τέταρτον, τρίτον δὲ πρὸς τέταρτον μείζονα λόγον ἔθῃ ἢ πέμπτον
πρὸς ἕκτον, καὶ πρῶτον πρὸς δεύτερον μείζονα λόγον ἕχει ἢ
πέμπτον πρὸς ἕκτον.}\\

\epsfysize=0.65in
\centerline{\epsffile{Book05/fig13g.eps}}

{\greekfont Πρῶτον γὰρ τὸ Α πρὸς δεύτερον τὸ Β τὸν α ὐτὸν ἐθέτω λόγον καὶ
τρίτον τὸ Γ πρὸς τέταρτον τὸ Δ, τρίτον δὲ τὸ Γ πρὸς τέταρτον τὸ Δ
μείζονα λόγον ἐθέτω ἢ πέμπτον τὸ Ε πρὸς ἕκτον τὸ Ζ. λέγω,
ὅτι καὶ πρῶτον τὸ Α πρὸς δεύτερον τὸ Β μείζονα λόγον ἕχει
ἤπερ πέμπτον τὸ Ε πρὸς ἕκτον τὸ Ζ.}

{\greekfont Ἐπεὶ γὰρ ἔστι τινὰ τῶν μὲν Γ, Ε ἰσάκις πολλαπλάσια, τῶν δὲ Δ, Ζ
ἄλλα, ἃ ἔτυθεν, ἰσάκις πολλαπλάσια, καὶ τὸ μὲν τοῦ Γ πολλαπλάσιον
τοῦ τοῦ Δ πολλαπλασίου ὑπερέθει, τὸ δὲ τοῦ Ε πολλαπλάσιον τοῦ τοῦ Ζ
πολλαπλασίου οὐθ ὑπερέθει, εἰλήφθω, καὶ ἔστω τῶν μὲν Γ, Ε ἰσάκις
πολλαπλάσια τὰ Η, Θ, τῶν δὲ Δ, Ζ ἄλλα, ἃ ἔτυθεν, ἰσάκις πολλαπλάσια
τὰ Κ, Λ, ὥστε τὸ μὲν Η τοῦ Κ ὑπερέθειν, τὸ δὲ Θ τοῦ Λ μὴ
ὑπερέθειν; καὶ ὁσαπλάσιον μέν ἐστι τὸ Η τοῦ Γ, τοσαυταπλάσιον
ἔστω καὶ τὸ Μ τοῦ Α, ὁσαπλάσιον δὲ τὸ Κ τοῦ Δ, τοσαυταπλάσιον
ἔστω καὶ τὸ Ν τοῦ Β.}

{\greekfont Καὶ ἐπεί ἐστιν ὡς τὸ Α πρὸς τὸ Β, οὕτως τὸ Γ πρὸς τὸ Δ, καὶ
εἴληπται τῶν μὲν Α, Γ ἰσάκις πολλαπλάσια τὰ Μ, Η, τῶν δὲ Β, Δ
ἄλλα, ἃ ἔτυθεν, ἰσάκις πολλαπλάσια τὰ Ν, Κ, εἰ ἄρα ὑπερέθει
τὸ Μ τοῦ Ν, ὑπερέθει καὶ τὸ Η τοῦ Κ, καὶ εἰ ἴσον, ἴσον,
καὶ εἰ ἔλαττον, ἔλλατον. ὑπερέθει δὲ τὸ Η τοῦ Κ; ὑπερέθει
ἄρα καὶ τὸ Μ τοῦ Ν. τὸ δὲ Θ τοῦ Λ οὐθ ὑπερέθει; καί
ἐστι τὰ μὲν Μ, Θ τῶν Α, Ε ἰσάκις πολλαπλάσια, τὰ δὲ Ν, Λ τῶν
Β, Ζ ἄλλα, ἃ ἔτυθεν, ἰσάκις πολλαπλάσια; τὸ ἄρα Α πρὸς
τὸ Β μείζονα λόγον ἔχει ἤπερ τὸ Ε πρὸς τὸ Ζ.}

{\greekfont Ἐὰν ἄρα πρῶτον πρὸς δεύτερον τὸν α ὐτὸν ἒθῃ  λόγον καὶ τρίτον
πρὸς τέταρτον, τρίτον δὲ πρὸς τέταρτον μείζονα λόγον ἔθῃ ἢ πέμπτον
πρὸς ἕκτον, καὶ πρῶτον πρὸς δεύτερον μείζονα λόγον ἕχει ἢ
πέμπτον πρὸς ἕκτον; ὅπερ ἔδει δεῖχαι.}}

\ParallelRText{
\begin{center}
{\large Proposition 13}$^\dag$
\end{center}

If a first (magnitude) has the same ratio to a second that a third (has) to a fourth, and the third (magnitude) has a greater ratio to the fourth than a fifth (has) to a sixth, (then) the first (magnitude) will also have a greater ratio to the
second  than the fifth (has) to the sixth.

\epsfysize=0.65in
\centerline{\epsffile{Book05/fig13e.eps}}

For let a first (magnitude) $A$ have the same ratio to a second $B$ that a
third $C$ (has) to a fourth $D$, and let the third (magnitude) $C$ have a
greater ratio to the fourth $D$  than a fifth $E$ (has) to a sixth $F$. I say that the first (magnitude) $A$
will also have a greater ratio to the second $B$ than the fifth $E$ (has) to the
sixth $F$.

For since there are some equal multiples of $C$ and $E$, and other random equal multiples of $D$ and $F$, (for which) the multiple of $C$ exceeds the (multiple) of $D$, 
and the multiple of $E$ does not exceed the multiple of $F$ [Def.~5.7], let them be taken. And let
$G$ and $H$ be equal multiples of $C$ and $E$ (respectively), and $K$ and $L$
other random equal multiples of $D$ and $F$ (respectively), such that 
$G$ exceeds $K$, but $H$ does not exceed $L$. And as many times as $G$ is (divisible)
by $C$, so many times let $M$ be (divisible) by $A$. And as many times as
$K$ (is divisible) by $D$, so many times let $N$ be (divisible) by $B$.

And since as $A$ is to $B$, so $C$ (is) to $D$, and the equal multiples $M$ and $G$
have been taken of $A$ and $C$ (respectively), and the other random equal
multiples $N$ and $K$ of $B$ and $D$ (respectively), thus if $M$ exceeds $N$ (then)
$G$ exceeds $K$, and if ($M$ is) equal (to $N$ then $G$ is also) equal (to $K$),
and if ($M$ is) less (than $N$ then $G$ is also) less (than $K$) [Def.~5.5]. And $G$ exceeds $K$. Thus, $M$ also exceeds $N$.
And $H$ does not exceed $L$. And $M$ and $H$ are equal multiples of
$A$ and $E$ (respectively), and $N$ and $L$ other random equal multiples of $B$ and $F$
(respectively).  Thus, $A$ has a greater ratio to $B$ than $E$ (has) to $F$  [Def.~5.7].

Thus, if a first (magnitude) has the same ratio to a second  that a third (has) to a fourth, and a third (magnitude) has a greater ratio to a fourth than a fifth (has) to a sixth, (then) the first (magnitude) will also have a greater ratio to the
second than the fifth (has) to the sixth. (Which is)
the very thing it was required to show.}
\end{Parallel}
{\footnotesize \noindent$^\dag$ In modern notation, this proposition
reads that if $\alpha:\beta::\gamma:\delta$ and $\gamma:\delta > \epsilon:\zeta$ then $\alpha:\beta>\epsilon:\zeta$.}

%%%%%%
% Prop 5.14
%%%%%%
\pdfbookmark[1]{Proposition 5.14}{pdf5.14}
\begin{Parallel}{}{} 
\ParallelLText{
\begin{center}
{\large \ggn{14}.}
\end{center}\vspace*{-7pt}

{\greekfont Ἐὰν πρῶτον πρὸς δεύτερον τὸν α ὐτὸν ἔθῃ λόγον καὶ τρίτον  πρὸς τέταρτον, τὸ δὲ
πρῶτον τοῦ τρίτου μεῖζον ᾖ, καὶ τὸ δεύτερον τοῦ τετάρτου
μεῖζον ἔσται, κἂν ἴσον, ἴσον, κἂν ἔλαττον, ἔλαττον.}\\~\\~\\

\epsfysize=0.5in
\centerline{\epsffile{Book05/fig14g.eps}}

{\greekfont Πρῶτον γὰρ τὸ Α πρὸς δεύτερον τὸ Β τὸν α ὐτὸν ἐθέτω λόγον καὶ
τρίτον τὸ Γ πρὸς τέταρτον τὸ Δ, μεῖζον δὲ ἔστω τὸ Α τοῦ Γ; λέγω,
ὅτι καὶ τὸ Β τοῦ Δ μεῖζόν ἐστιν.}

{\greekfont Ἐπεὶ γὰρ τὸ Α τοῦ Γ μεῖζόν ἐστιν, ἄλλο δέ, ὃ ἔτυθεν,
[μέγεθος] τὸ Β, τὸ Α ἄρα πρὸς τὸ Β μείζονα λόγον ἔχει
ἤπερ τὸ Γ πρὸς τὸ Β. ὡς δὲ τὸ Α 	πρὸς τὸ Β, οὕτως τὸ Γ πρὸς
τὸ
Δ; καὶ τὸ Γ ἄρα πρὸς τὸ Δ μείζονα
λόγον ἔχει ἤπερ τὸ Γ πρὸς τὸ Β. πρὸς ὃ δὲ τὸ α ὐτὸ μείζονα
λόγον ἔχει, ἐκεῖνο ἔλασσόν ἐστιν; ἔλασσον ἄρα τὸ Δ τοῦ Β; ὥστε
μεῖζόν ἐστι τὸ Β τοῦ Δ.}

{\greekfont Ὁμοίως δὴ δεῖχομεν, ὅτι κἂν ἴσον ᾖ τὸ Α τῷ Γ,
ἴσον ἔσται καὶ τὸ Β τῷ Δ, κἄν ἔλασσον ᾖ τὸ Α
τοῦ Γ, ἔλασσον ἔσται καὶ τὸ Β τοῦ Δ.}

{\greekfont Ἐὰν ἄρα πρῶτον πρὸς δεύτερον τὸν α ὐτὸν ἔθῃ λόγον καὶ τρίτον  πρὸς τέταρτον, τὸ δὲ
πρῶτον τοῦ τρίτου μεῖζον ᾖ, καὶ τὸ δεύτερον τοῦ τετάρτου
μεῖζον ἔσται, κἂν ἴσον, ἴσον, κἂν ἔλαττον, ἔλαττον;
ὅπερ ἔδει δεῖχαι.
}}

\ParallelRText{
\begin{center}
{\large Proposition 14}$^\dag$
\end{center}

If a first (magnitude) has the same ratio to a second  that a third (has) to a fourth, and the first (magnitude) is greater than the third,
(then) the second will also be greater than the fourth. And if (the first magnitude is) equal
(to the third then the second will also be) equal (to the fourth). And
if (the first magnitude is) less (than the third then the second will also be) less (than the fourth).

\epsfysize=0.5in
\centerline{\epsffile{Book05/fig14e.eps}}

For let a first (magnitude) $A$ have the same ratio to a second $B$  that a
third $C$ (has) to a fourth $D$. And let $A$ be greater than $C$. I say that $B$
is also greater than $D$.

For since $A$ is greater than $C$, and $B$ (is) another random [magnitude],
$A$ thus has a greater ratio to $B$ than $C$ (has) to $B$ [Prop.~5.8]. And as $A$ (is) to $B$, so $C$ (is) to
$D$. Thus, $C$ also has a greater ratio to $D$ than $C$ (has) to $B$. And that (magnitude) to
which the same (magnitude) has a greater ratio is the lesser
 [Prop.~5.10]. Thus, $D$
(is) less than $B$. Hence, $B$ is greater than $D$.

So, similarly, we can show that even if $A$ is equal to $C$, (then) $B$ will also be equal to $D$, and even if $A$ is less than $C$, (then) $B$ will also be less than $D$.

Thus, if a first (magnitude) has the same ratio to a second  that a third (has) to a fourth, and the first (magnitude) is greater than the third,
(then) the second will also be greater than the fourth. And if (the first magnitude is) equal
(to the third then the second will also be) equal (to the fourth). And
if (the first magnitude  is) less (than the third then the second will also be) less (than the fourth). (Which is) the very thing it was required to show.}
\end{Parallel}
{\footnotesize \noindent$^\dag$ In modern notation, this proposition
reads that if $\alpha:\beta::\gamma:\delta$ then $\alpha \gtreqqless \gamma$ as
$\beta\gtreqqless\delta$.}

%%%%%%
% Prop 5.15
%%%%%%
\pdfbookmark[1]{Proposition 5.15}{pdf5.15}
\begin{Parallel}{}{} 
\ParallelLText{
\begin{center}
{\large \ggn{15}.}
\end{center}\vspace*{-7pt}

{\greekfont Τὰ μέρη τοῖς ὡσαύτως πολλαπλασίοις τὸν α ὐτὸν ἔχει λόγον
ληφθέντα κατάλληλα.}

\epsfysize=0.8in
\centerline{\epsffile{Book05/fig15g.eps}}

{\greekfont Ἔστω γὰρ ἰσάκις πολλαπλάσιον τὸ ΑΒ τοῦ Γ καὶ το ΔΕ τοῦ Ζ;
λέγω, ὅτι ἐστὶν ὡς τὸ Γ πρὸς τὸ Ζ, οὕτως τὸ ΑΒ πρὸς τὸ ΔΕ.}

{\greekfont Ἐπεὶ γὰρ ἰσάκις ἐστὶ πολλαπλάσιον τὸ ΑΒ τοῦ Γ καὶ τὸ ΔΕ τοῦ Ζ,
ὅσα ἄρα ἐστὶν ἐν τῷ ΑΒ μεγέθη ἴσα τῷ Γ, τοσαῦτα καὶ
ἐν τῷ ΔΕ ἴσα τῷ Ζ. διῃρήσθω τὸ μὲν ΑΒ εἰς τὰ τῷ Γ ἴσα
τὰ ΑΗ, ΗΘ, ΘΒ, τὸ δὲ ΔΕ εἰς τὰ τῷ Ζ ἴσα τὰ ΔΚ, ΚΛ, ΛΕ;
ἔσται δὴ ἴσον τὸ πλῆθος τῶν ΑΗ, ΗΘ, ΘΒ τῷ πλήθει τῶν ΔΚ, ΚΛ, ΛΕ. καὶ ἐπεὶ ἴσα ἐστὶ τὰ ΑΗ, ΗΘ, ΘΒ  ἀλλήλοις, ἔστι δὲ καὶ τὰ
ΔΚ, ΚΛ, ΛΕ ἴσα ἀλλήλοις,
ἔστιν ἄρα
ὡς τὸ ΑΗ πρὸς τὸ ΔΚ,  οὕτως τὸ ΗΘ πρὸς τὸ ΚΛ, καὶ τὸ ΘΒ
πρὸς τὸ ΛΕ. ἔσται ἄρα καὶ ὡς ἓν τῶν ἡγουμένων 
πρὸς ἓν τῶν ἑπομένων, οὕτως ἅπαντα τὰ ἡγουμένα πρὸς
ἅπαντα τὰ ἑπόμενα; ἔστιν ἄρα ὡς τὸ ΑΗ πρὸς τὸ ΔΚ, οὕτως
τὸ ΑΒ πρὸς τὸ ΔΕ. ἴσον δὲ τὸ μὲν ΑΗ τῷ Γ, τὸ δὲ ΔΚ
τῷ Ζ; ἔστιν ἄρα ὡς τὸ Γ πρὸς τὸ Ζ οὕτως τὸ ΑΒ πρὸς 
τὸ ΔΕ.}

{\greekfont Τὰ ἄρα μέρη τοῖς ὡσαύτως πολλαπλασίοις τὸν α ὐτὸν ἔχει λόγον
ληφθέντα κατάλληλα; ὅπερ ἔδει δεῖχαι.}}

\ParallelRText{
\begin{center}
{\large Proposition 15}$^\dag$
\end{center}

Parts have the same ratio  as similar multiples,
taken in corresponding order.

\epsfysize=0.8in
\centerline{\epsffile{Book05/fig15e.eps}}

For let $AB$ and $DE$ be equal multiples of $C$ and $F$ (respectively). I say that
as $C$ is to $F$, so $AB$ (is) to $DE$.

For since $AB$ and $DE$ are equal multiples of $C$ and $F$ (respectively), thus
as many magnitudes as there are in $AB$ equal to $C$, so many (are there) also
in $DE$ equal to $F$. Let $AB$ be divided into (magnitudes) $AG$, $GH$, $HB$, equal to $C$, and $DE$ into (magnitudes) $DK$, $KL$,  $LE$, equal to
$F$. So, the number of (magnitudes) $AG$, $GH$,  $HB$ will equal the number
of (magnitudes) $DK$, $KL$, $LE$.
And since $AG$, $GH$, $HB$ are equal to one another, 
and $DK$, $KL$,  $LE$ are also equal to one another,
thus as $AG$ is to
$DK$, so $GH$ (is) to $KL$, and $HB$ to $LE$ [Prop.~5.7]. And, thus (for proportional magnitudes), as one of the leading
(magnitudes) will be to one of the following, so all of the leading (magnitudes will be)
to all of the following [Prop.~5.12].
Thus, as $AG$ is to $DK$, so $AB$ (is) to $DE$. And $AG$ is equal to $C$, and 
$DK$ to $F$. Thus, as $C$ is to $F$, so $AB$ (is) to $DE$.

Thus, parts have the same ratio  as similar multiples,
taken in corresponding order. (Which is) the very thing it was required to
show.}
\end{Parallel}
{\footnotesize \noindent$^\dag$ In modern notation, this proposition
reads that $\alpha:\beta::m\,\alpha:m\,\beta$.}

%%%%%%
% Prop 5.16
%%%%%%
\pdfbookmark[1]{Proposition 5.16}{pdf5.16}
\begin{Parallel}{}{} 
\ParallelLText{
\begin{center}
{\large \ggn{16}.}
\end{center}\vspace*{-7pt}

{\greekfont Ἐὰν τέσσαρα μεγέθη ἀνάλογον ᾖ, καὶ ἐναλλὰχ ἀνάλογον ἔσται.}

{\greekfont Ἔστω τέσσαρα μεγέθη ἀνάλογον τὰ 
Α, Β, Γ, Δ, ὡς τὸ Α πρὸς τὸ Β, οὕτως τὸ Γ πρὸς τὸ Δ; λέγω,
ὅτι καὶ ἐναλλὰχ [ἀνάλογον] ἔσται, ὡς τὸ Α πρὸς τὸ Γ, 
οὕτως τὸ Β πρὸς τὸ Δ.}

{\greekfont Εἰλήφθω γὰρ τῶν μὲν Α, Β ἰσάκις πολλαπλάσια τὰ Ε, Ζ, τῶν
δὲ Γ, Δ ἄλλα, ἃ ἔτυθεν, ἰσάκις πολλαπλάσια τὰ Η, Θ.}\\

\epsfysize=1in
\centerline{\epsffile{Book05/fig16g.eps}}

{\greekfont Καὶ ἐπεὶ ἰσάκις ἐστὶ πολλαπλάσιον τὸ Ε τοῦ Α καὶ τὸ Ζ τοῦ
Β, τὰ δὲ μέρη τοῖς ὡσαύτως πολλαπλασίοις τὸν α ὐτὸν ἔχει
λόγον, ἔστιν ἄρα ὡς τὸ Α πρὸς
τὸ Β, οὕτως τὸ Ε πρὸς τὸ Ζ. ὡς δὲ τὸ Α πρὸς τὸ Β,
οὕτως τὸ Γ πρὸς τὸ Δ; καὶ ὡς ἄρα τὸ Γ πρὸς
τὸ Δ, οὕτως τὸ Ε πρὸς τὸ Ζ. πάλιν, ἐπεὶ τὰ Η, Θ τῶν Γ, Δ
ἰσάκις ἐστὶ πολλαπλάσια, ἔστιν ἄρα ὡς τὸ Γ πρὸς τὸ Δ, οὕτως
τὸ Η πρὸς τὸ Θ. ὡς δὲ τὸ Γ πρὸς τὸ Δ, [οὕτως] τὸ Ε πρὸς
τὸ Ζ; καὶ ὡς ἄρα τὸ Ε πρὸς τὸ Ζ, οὕτως τὸ Η πρὸς τὸ Θ. ἐὰν
δὲ τέσσαρα μεγέθη ἀνάλογον ᾖ, τὸ δὲ πρῶτον τοῦ τρίτου
μεῖζον ᾖ, καὶ τὸ δεύτερον τοῦ τετάρτου μεῖζον ἔσται,
κἂν ἴσον, ἴσον, κἄν ἔλαττον, ἔλαττον. εἰ ἄρα ὑπερέθει
τὸ Ε τοῦ Η, ὑπερέθει καὶ τὸ Ζ τοῦ Θ, καὶ εἰ ἴσον, ἴσον,
καὶ εἰ ἔλαττον, ἔλαττον. καί ἐστι τὰ μὲν Ε, Ζ τῶν Α, Β
ἰσάκις πολλαπλάσια, τὰ δὲ Η, Θ τῶν Γ, Δ ἄλλα, ἃ ἔτυθεν, ἰσάκις
πολλαπλάσια; ἔστιν ἄρα ὡς τὸ Α πρὸς τὸ Γ, οὕτως τὸ Β πρὸς
τὸ Δ.}

{\greekfont Ἐὰν ἄρα τέσσαρα μεγέθη ἀνάλογον ᾖ, καὶ ἐναλλὰχ
ἀνάλογον ἔσται; ὅπερ ἔδει δεῖχαι.}}

\ParallelRText{
\begin{center}
{\large Proposition 16}$^\dag$
\end{center}

If four magnitudes are proportional, (then) they will
also be proportional alternately.

Let $A$, $B$, $C$ and $D$ be four proportional magnitudes, (such that) as $A$
(is) to $B$, so $C$ (is) to $D$. I say that they will also be [proportional]
 alternately, (so that) as $A$ (is) to $C$, so $B$ (is) to $D$.
 
For let the equal multiples $E$ and $F$ be taken of $A$ and $B$ (respectively),
and the other random equal multiples $G$ and $H$ of $C$ and $D$ (respectively).

\epsfysize=1in
\centerline{\epsffile{Book05/fig16e.eps}}

And since $E$ and $F$ are equal multiples of $A$ and $B$ (respectively),
and parts  have the same ratio as similar multiples [Prop.~5.15], thus as $A$ is to $B$, so $E$ (is) to $F$.
But as $A$ (is) to $B$, so $C$ (is) to $D$. And, thus, as $C$ (is) to $D$, so
$E$ (is) to $F$ [Prop.~5.11].
Again, since $G$ and $H$ are equal multiples of $C$ and $D$ (respectively),  thus
as $C$  is to $D$, so $G$ (is) to $H$ [Prop.~5.15].
But as $C$ (is) to $D$, [so] $E$ (is) to $F$. And, thus, as $E$ (is) to $F$, so
$G$ (is) to $H$ [Prop.~5.11]. And if four magnitudes are proportional, and the first is greater than the third, (then)
the second will also be greater than the fourth, and if (the first is) equal
(to the third then the second will also be) equal (to the fourth), and if
(the first is) less (than the third then the second will also be) less (than the fourth)
[Prop.~5.14]. Thus, if $E$ exceeds $G$ (then)
$F$ also exceeds $H$, and if ($E$ is) equal (to $G$ then $F$ is also) equal (to $H$), and if
($E$ is) less (than $G$ then $F$ is also) less (than $H$). And $E$ and $F$ are equal multiples of $A$ and $B$ (respectively), and $G$ and $H$ other random equal multiples
of $C$ and $D$ (respectively). Thus, as $A$ is to $C$, so $B$ (is) to $D$ [Def.~5.5].

Thus, if four magnitudes are proportional, (then) they will
also be proportional  alternately. (Which is) the very thing it was required to show.}
\end{Parallel}
{\footnotesize \noindent$^\dag$ In modern notation, this proposition
reads that if $\alpha:\beta::\gamma:\delta$ then $\alpha:\gamma::\beta:\delta$.}

%%%%%%
% Prop 5.17
%%%%%%
\pdfbookmark[1]{Proposition 5.17}{pdf5.17}
\begin{Parallel}{}{} 
\ParallelLText{
\begin{center}
{\large \ggn{17}.}
\end{center}\vspace*{-7pt}

{\greekfont Ἐὰν συγκείμενα μεγέθη ἀνάλογον ᾖ, καὶ διαιρεθέντα ἀνάλογον
ἔσται.}

\epsfysize=1in
\centerline{\epsffile{Book05/fig17g.eps}}

{\greekfont Ἔστω συγκείμενα μεγέθη ἀνάλογον τὰ ΑΒ, ΒΕ, ΓΔ, ΔΖ, ὡς τὸ ΑΒ
πρὸς τὸ ΒΕ, οὕτως τὸ ΓΔ πρὸς τὸ ΔΖ; λέγω, ὅτι καὶ διαιρεθέντα
ἀνάλογον ἔσται, ὡς τὸ ΑΕ πρὸς τὸ ΕΒ, οὕτως τὸ ΓΖ πρὸς τὸ ΔΖ.}

{\greekfont Εἰλήφθω γὰρ τῶν μὲν ΑΕ, ΕΒ, ΓΖ, ΖΔ ἰσάκις πολλαπλάσια τὰ ΗΘ,
ΘΚ, ΛΜ, ΜΝ, τῶν δὲ ΕΒ, ΖΔ ἄλλα, ἃ ἔτυθεν, ἰσάκις
πολλαπλάσια τὰ ΚΧ, ΝΠ.}

{\greekfont Καὶ ἐπεὶ ἰσάκις ἐστὶ πολλαπλάσιον τὸ ΗΘ τοῦ ΑΕ καὶ τὸ ΘΚ
τοῦ ΕΒ, ἰσάκις ἄρα ἐστὶ πολλαπλάσιον τὸ ΗΘ τοῦ ΑΕ καὶ τὸ ΗΚ
τοῦ ΑΒ. ἰσάκις δέ ἐστι πολλαπλάσιον τὸ ΗΘ τοῦ ΑΕ καὶ τὸ ΛΜ
τοῦ ΓΖ; ἰσάκις ἄρα ἐστὶ πολλαπλάσιον τὸ ΗΚ τοῦ ΑΒ καὶ τὸ ΛΜ
τοῦ ΓΖ. πάλιν, ἐπεὶ ἰσάκις ἐστὶ πολλαπλάσιον τὸ ΛΜ τοῦ ΓΖ
καὶ τὸ ΜΝ τοῦ ΖΔ, ἰσάκις ἄρα ἐστὶ πολλαπλάσιον τὸ ΛΜ τοῦ ΓΖ
καὶ τὸ ΛΝ τοῦ ΓΔ. ἰσάκις δὲ ἦν πολλαπλάσιον τὸ ΛΜ τοῦ ΓΖ καὶ τὸ
ΗΚ τοῦ ΑΒ; ἰσάκις ἄρα ἐστὶ πολλαπλάσιον τὸ ΗΚ τοῦ ΑΒ καὶ τὸ ΛΝ
τοῦ ΓΔ. τὰ ΗΚ, ΛΝ ἄρα τῶν ΑΒ, ΓΔ ἰσάκις ἐστὶ πολλαπλάσια.
πάλιν, ἐπεὶ ἰσάκις ἐστὶ πολλαπλασίον τὸ ΘΚ τοῦ ΕΒ
καὶ τὸ ΜΝ τοῦ ΖΔ, ἔστι δὲ καὶ τὸ ΚΧ τοῦ ΕΒ ἰσάκις πολλαπλάσιον
καὶ τὸ ΝΠ τοῦ ΖΔ, καὶ συντεθὲν τὸ ΘΧ τοῦ ΕΒ ἰσάκις ἐστὶ
πολλαπλάσιον καὶ τὸ ΜΠ τοῦ ΖΔ. 
καὶ ἐπεί ἐστιν ὡς τὸ ΑΒ πρὸς  τὸ ΒΕ, οὕτως τὸ ΓΔ πρὸς τὸ ΔΖ, καὶ εἴληπται τῶν μὲν ΑΒ, ΓΔ ἰσάκις πολλαπλάσια τὰ
ΗΚ, ΛΝ, τῶν δὲ ΕΒ, ΖΔ ἰσάκις πολλαπλάσια τὰ ΘΧ, ΜΠ, εἰ ἄρα
ὑπερέθει τὸ ΗΚ τοῦ ΘΧ, ὑπερέθει καὶ τὸ ΛΝ τοῦ ΜΠ, καὶ εἰ
ἴσον, ἴσον, καὶ εἰ ἔλαττον, ἔλαττον. ὑπερεθέτω δὴ τὸ ΗΚ τοῦ
ΘΧ, καὶ κοινοῦ ἀφαιρεθέντος τοῦ ΘΚ ὑπερέθει ἄρα καὶ τὸ
ΗΘ τοῦ ΚΧ. ἀλλα εἰ ὑπερεῖθε
τὸ ΗΚ τοῦ ΘΧ
ὑπερεῖθε καὶ τὸ ΛΝ τοῦ ΜΠ; ὑπερέθει ἄρα καὶ τὸ ΛΝ τοῦ ΜΠ,
καὶ κοινοῦ ἀφαιρεθέντος τοῦ ΜΝ ὑπερέθει
καὶ τὸ ΛΜ τοῦ ΝΠ; ὥστε εἰ ὑπερέθει τὸ ΗΘ τοῦ ΚΧ,
ὑπερέθει καὶ τὸ ΛΜ τοῦ  ΝΠ. ὁμοίως δὴ δεῖχομεν, ὅτι κἂν
ἴσον ᾖ τὸ ΗΘ τῷ ΚΧ, ἴσον ἔσται καὶ τὸ ΛΜ τῷ ΝΠ, κἂν
ἔλαττον, ἔλαττον. καί ἐστι τὰ μὲν ΗΘ, ΛΜ τῶν ΑΕ, ΓΖ
ἰσάκις πολλαπλάσια, τὰ δὲ ΚΧ, ΝΠ τῶν ΕΒ, ΖΔ ἄλλα, ἃ ἔτυθεν,
ἰσάκις πολλαπλάσια; ἔστιν ἄρα ὡς τὸ ΑΕ πρὸς τὸ ΕΒ, οὕτως τὸ
ΓΖ πρὸς τὸ ΖΔ.}

{\greekfont Ἐὰν ἄρα συγκείμενα μεγέθη ἀνάλογον ᾖ, καὶ διαιρεθέντα ἀνάλογον
ἔσται; ὅπερ ἔδει δεῖχαι.}}

\ParallelRText{
\begin{center}
{\large Proposition 17}$^\dag$
\end{center}

If composed  magnitudes are proportional, (then) they will also be proportional (when)
separarted.

\epsfysize=1in
\centerline{\epsffile{Book05/fig17e.eps}}

Let $AB$, $BE$, $CD$, and $DF$ be composed magnitudes (which are) proportional,
(so that) as $AB$ (is) to $BE$, so $CD$ (is) to $DF$. I say that they will also be
proportional (when) separated, (so that) as $AE$ (is) to $EB$, so  $CF$ (is)
to $DF$.

For let the equal multiples $GH$, $HK$, $LM$, and $MN$ be taken of
$AE$, $EB$, $CF$, and $FD$ (respectively), and the other random equal multiples
$KO$ and $NP$ of $EB$ and $FD$ (respectively).

And since $GH$ and $HK$ are equal multiples of $AE$ and $EB$ (respectively), 
$GH$ and $GK$ are thus equal multiples of $AE$ and $AB$ (respectively) [Prop.~5.1]. But $GH$ and $LM$ are equal
multiples of $AE$ and $CF$ (respectively). Thus, $GK$ and $LM$ are equal
multiples of $AB$ and $CF$ (respectively). Again, since $LM$ and $MN$ are equal
multiples of $CF$ and $FD$ (respectively), $LM$ and $LN$ are thus equal multiples
of $CF$ and $CD$ (respectively) [Prop.~5.1]. 
And $LM$ and $GK$ were equal multiples of $CF$ and $AB$ (respectively).
Thus, $GK$ and $LN$ are equal multiples of $AB$ and $CD$ (respectively).
Thus, $GK$, $LN$ are equal multiples of $AB$, $CD$. Again, since
$HK$ and $MN$ are equal multiples of $EB$ and $FD$ (respectively), and $KO$
and $NP$ are also  equal multiples of $EB$ and $FD$ (respectively),  (then), added together, $HO$ and $MP$ are also equal multiples of $EB$ and $FD$ (respectively) [Prop.~5.2]. And since as $AB$ (is) to $BE$, so
$CD$ (is) to $DF$, and the equal multiples $GK$, $LN$ be
taken of $AB$, $CD$, and the equal multiples $HO$, $MP$ of $EB$, $FD$, thus 
if $GK$ exceeds $HO$ (then) $LN$ also exceeds $MP$, and if ($GK$ is) equal
(to $HO$ then $LN$ is also) equal (to $MP$), and if ($GK$ is) less (than $HO$ then
$LN$ is also) less (than $MP$) [Def.~5.5]. So let $GK$ exceed $HO$, and thus, $HK$ being taken away from
both, $GH$ exceeds $KO$. But (we saw that) if $GK$ was exceeding $HO$ (then)
$LN$ was also exceeding $MP$. Thus, $LN$ also exceeds $MP$, and, $MN$ being
taken away from both, $LM$ also exceeds $NP$.  Hence, if $GH$ exceeds $KO$
(then) $LM$ also exceeds $NP$. So, similarly, we can show that even if
$GH$ is equal to $KO$ (then) $LM$ will also be equal to $NP$, and
even if ($GH$ is) less (than $KO$ then $LM$ will also be) less (than $NP$).
And $GH$, $LM$ are equal multiples of $AE$, $CF$, and $KO$, $NP$ other random
equal multiples of $EB$, $FD$. Thus, as $AE$ is to $EB$, so  $CF$ (is) to $FD$ [Def.~5.5].

Thus, if composed magnitudes are proportional, (then) they will also be proportional (when)
separarted. (Which is) the very thing it was required to show.}
\end{Parallel}
{\footnotesize \noindent$^\dag$ In modern notation, this proposition
reads that if $\alpha+\beta:\beta::\gamma+\delta:\delta$ then $\alpha:\beta::
\gamma:\delta$.}

%%%%%%
% Prop 5.18
%%%%%%
\pdfbookmark[1]{Proposition 5.18}{pdf5.18}
\begin{Parallel}{}{} 
\ParallelLText{
\begin{center}
{\large \ggn{18}.}
\end{center}\vspace*{-7pt}

{\greekfont Ἐὰν διῃρημένα μεγέθη ἀνάλογον ᾖ, καὶ συντεθέντα ἀνάλογον
ἔσται.}

\epsfysize=0.7in
\centerline{\epsffile{Book05/fig18g.eps}}

{\greekfont Ἔστω διῃρημένα μεγέθη ἀνάλογον τὰ ΑΕ, ΕΒ, ΓΖ, ΖΔ, ὡς τὸ ΑΕ
πρὸς τὸ ΕΒ, οὕτως τὸ ΓΖ πρὸς τὸ ΖΔ; λέγω, ὅτι καὶ συντεθέντα
ἀνάλογον ἔσται, ὡς τὸ ΑΒ πρὸς τὸ ΒΕ, οὕτως τὸ ΓΔ πρὸς τὸ ΖΔ.}

{\greekfont Εἰ γὰρ μή ἐστὶν ὡς τὸ ΑΒ πρὸς τὸ ΒΕ, οὕτως τὸ ΓΔ πρὸς
τὸ ΔΖ, ἔσται ὡς τὸ ΑΒ πρὸς τὸ ΒΕ, οὕτως τὸ ΓΔ ἤτοι πρὸς ἔλασσόν τι τοῦ ΔΖ ἢ πρὸς μεῖζον.}

{\greekfont Ἔστω πρότερον πρὸς ἔλασσον τὸ ΔΗ. καὶ ἐπεί ἐστιν ὡς τὸ ΑΒ
πρὸς τὸ ΒΕ, οὕτως τὸ ΓΔ πρὸς τὸ ΔΗ,  συγκείμενα μεγέθη ἀνάλογόν
ἐστιν; ὥστε καὶ διαιρεθέντα ἀνάλογον ἔσται. ἔστιν ἄρα ὡς τὸ ΑΕ
πρὸς τὸ ΕΒ, οὕτως τὸ ΓΗ πρὸς τὸ ΗΔ. ὑπόκειται δὲ καὶ ὡς τὸ
ΑΕ πρὸς τὸ ΕΒ, οὕτως τὸ ΓΖ πρὸς τὸ ΖΔ. καὶ ὡς ἄρα τὸ ΓΗ πρὸς
τὸ ΗΔ, οὕτως τὸ ΓΖ πρὸς τὸ ΖΔ. μεῖζον δὲ τὸ πρῶτον τὸ ΓΗ τοῦ
τρίτου τοῦ ΓΖ; μεῖζον ἄρα καὶ τὸ δεύτερον
τὸ ΗΔ τοῦ τετάρτου τοῦ ΖΔ. ἀλλὰ καὶ ἔλαττον; ὅπερ
ἐστὶν ἀδύνατον; οὐκ ἄρα ἐστὶν ὡς τὸ ΑΒ πρὸς τὸ ΒΕ, οὕτως
τὸ ΓΔ πρὸς ἔλασσον
 τοῦ ΖΔ. ὁμοίως δὴ δείχομεν, ὅτι
οὐδὲ πρὸς μεῖζον; πρὸς α ὐτὸ ἄρα.}

{\greekfont Ἐὰν ἄρα διῃρημένα μεγέθη ἀνάλογον ᾖ, καὶ συντεθέντα ἀνάλογον
ἔσται; ὅπερ ἔδει δεῖχαι.}}

\ParallelRText{
\begin{center}
{\large Proposition 18}$^\dag$
\end{center}

If separated  magnitudes are proportional,  (then) they will also be proportional (when)
composed.

\epsfysize=0.7in
\centerline{\epsffile{Book05/fig18e.eps}}

Let $AE$, $EB$, $CF$, and $FD$ be separated magnitudes (which are) proportional,
(so that) as $AE$ (is) to $EB$, so  $CF$ (is) to $FD$. I say that they will also be proportional (when)  composed, (so that) as $AB$ (is) to $BE$, so  $CD$ (is) to $FD$.

For if (it is) not (the case that) as $AB$ is to $BE$, so $CD$ (is) to $FD$, (then) it will
surely be (the case that) as $AB$ (is) to $BE$, so
 $CD$ is either to some (magnitude) less than $DF$, or (some magnitude) greater (than $DF$).$^\ddag$
 
 Let it, first of all, be to (some magnitude) less (than $DF$), (namely) $DG$. And since composed magnitudes are proportional, (so that) as $AB$ is
 to $BE$, so $CD$ (is) to $DG$,  they will  thus also be proportional (when) separated 
 [Prop.~5.17]. Thus, as $AE$ is to $EB$, so
 $CG$ (is) to $GD$. But it was also assumed that as $AE$ (is) to $EB$, so $CF$ (is)
 to $FD$. Thus, (it is) also (the case that) as $CG$ (is) to $GD$, so $CF$ (is) to $FD$ [Prop.~5.11]. And the first (magnitude) $CG$ (is)
 greater than the third $CF$. Thus, the second (magnitude) $GD$ (is)
 also greater than the fourth $FD$ [Prop.~5.14]. 
 But (it is) also less. The very thing is impossible. Thus, (it is) not (the case that) as $AB$ is to $BE$, so $CD$ (is) to less than $FD$. Similarly, we can show
 that neither (is it the case) to greater (than $FD$). Thus, (it is the case) to
 the same (as $FD$).
 
 Thus, if separated  magnitudes are proportional,  (then) they will also be 
 proportional (when)
composed. (Which is) the very thing it was required to show.}
\end{Parallel}
{\footnotesize \noindent$^\dag$ In modern notation, this proposition
reads that if $\alpha:\beta::\gamma:\delta$ then $\alpha+\beta:\beta::
\gamma+\delta:\delta$.}\\
{\footnotesize \noindent$^\ddag$ Here, Euclid assumes, without proof, that
a fourth magnitude proportional to three given magnitudes can always
be found.}

%%%%%%
% Prop 5.19
%%%%%%
\pdfbookmark[1]{Proposition 5.19}{pdf5.19}
\begin{Parallel}{}{} 
\ParallelLText{
\begin{center}
{\large \ggn{19}.}
\end{center}\vspace*{-7pt}

{\greekfont Ἐὰν ᾖ ὡς ὅλον πρὸς ὅλον, οὕτως ἀφαιρεθὲν πρὸς
ἀφαιρεθέν, καὶ τὸ λοιπὸν πρὸς τὸ λοιπὸν ἔσται ὡς ὅλον
πρὸς ὅλον.}\\

\epsfysize=0.6in
\centerline{\epsffile{Book05/fig19g.eps}}

{\greekfont Ἔστω γὰρ ὡς ὅλον τὸ ΑΒ πρὸς ὅλον τὸ ΓΔ, οὕτως ἀφαιρεθὲν
τὸ ΑΕ πρὸς ἀφαιρεθὲν τὸ ΓΖ; λέγω, ὅτι καὶ λοιπὸν τὸ ΕΒ πρὸς
λοιπὸν τὸ ΖΔ ἔσται ὡς ὅλον τὸ ΑΒ πρὸς ὅλον τὸ ΓΔ.}

{\greekfont Ἐπεὶ γάρ ἐστιν ὡς τὸ ΑΒ πρὸς τὸ ΓΔ, οὕτως τὸ ΑΕ πρὸς τὸ
ΓΖ, καὶ ἐναλλὰχ ὡς τὸ ΒΑ πρὸς τὸ ΑΕ, οὕτως τὸ ΔΓ πρὸς τὸ ΓΖ.
καὶ ἐπεὶ συγκείμενα μεγέθη ἀνάλογόν ἐστιν, καὶ διαιρεθέντα
ἀνάλογον ἔσται, ὡς τὸ ΒΕ πρὸς τὸ ΕΑ, οὕτως τὸ ΔΖ πρὸς
τὸ ΓΖ; καὶ ἐναλλάχ, ὡς τὸ ΒΕ πρὸς τὸ ΔΖ, οὕτως τὸ ΕΑ πρὸς
τὸ ΖΓ. ὡς δὲ τὸ ΑΕ πρὸς τὸ ΓΖ,  οὕτως ὑπόκειται ὅλον τὸ ΑΒ πρὸς ὅλον
τὸ ΓΔ. καὶ λοιπὸν ἄρα τὸ ΕΒ πρὸς λοιπὸν τὸ ΖΔ ἔσται ὡς ὅλον
τὸ ΑΒ πρὸς ὅλον τὸ ΓΔ.}

{\greekfont Ἐὰν ἄρα ᾖ ὡς ὅλον πρὸς ὅλον, οὕτως ἀφαιρεθὲν πρὸς
ἀφαιρεθέν, καὶ τὸ λοιπὸν πρὸς τὸ λοιπὸν ἔσται ὡς ὅλον
πρὸς ὅλον [ὅπερ ἔδει δεῖχαι].}

{\greekfont μ̀βοχ{[}Καὶ ἐπεὶ ἐδείθθη ὡς τὸ ΑΒ πρὸς τὸ ΓΔ, οὕτως τὸ ΕΒ
πρὸς τὸ ΖΔ, καὶ ἐναλλὰχ ὡς τὸ ΑΒ πρὸς τὸ ΒΕ οὕτως τὸ
ΓΔ πρὸς τὸ ΖΔ, συγκείμενα ἄρα μεγέθη ἀνάλογόν ἐστιν;
ἐδείθθη δὲ ὡς τὸ ΒΑ πρὸς τὸ ΑΕ, οὕτως τὸ ΔΓ πρὸς τὸ ΓΖ;
καί ἐστιν ἀναστρέψαντι].}\\~\\

\begin{center}
{\large {\greekfont Πόρισμα}.}
\end{center}\vspace*{-7pt}

{\greekfont Ἐκ δὴ τούτου φανερόν, ὅτι ἐὰν συγκείμενα μεγέθη
ἀνάλογον ᾖ, καὶ ἀναστρέψαντι ἀνάλογον ἔσται; ὅπερ ἔδει
δεῖχαι.}}

\ParallelRText{
\begin{center}
{\large Proposition 19}$^\dag$
\end{center}

If as the whole is to the whole so the (part) taken
away is to the (part) taken away, (then) the remainder to the remainder will
also be as the whole (is) to the whole.

\epsfysize=0.6in
\centerline{\epsffile{Book05/fig19e.eps}}

For let the whole $AB$ be to the whole $CD$ as the (part) taken away $AE$
(is) to the (part) taken away $CF$. I say that the remainder $EB$
to the remainder $FD$ will also be as the whole $AB$ (is) to the whole $CD$.

For since as $AB$ is to $CD$, so $AE$ (is) to $CF$, (it is) also (the case), alternately, 
(that) as $BA$ (is)
to $AE$, so $DC$ (is) to $CF$ [Prop.~5.16].
And since composed magnitudes are proportional, (then) they will also
be proportional (when) separated, (so that) as $BE$ (is) to $EA$, so $DF$ (is)
to $CF$ [Prop.~5.17].
Also, alternately, as $BE$ (is) to $DF$, so $EA$ (is) to $FC$ [Prop.~5.16]. And it was assumed that as
$AE$ (is) to $CF$, so the whole $AB$ (is) to the whole $CD$. And, thus,  as the
remainder $EB$ (is) to the remainder $FD$, so the whole $AB$ will be to the
whole $CD$.

Thus, if as the whole is to the whole so the (part) taken
away is to the (part) taken away, (then) the remainder to the remainder will
also be as the whole (is) to the whole. [(Which is) the very thing it was required to show.]

\mbox{[}And since it was shown (that) as $AB$ (is) to $CD$, so $EB$ (is) to $FD$, (it is)
also (the case), alternately, (that) as $AB$ (is) to $BE$, so $CD$ (is) to $FD$. 
Thus, composed magnitudes are proportional. And it was shown
(that) as $BA$ (is) to $AE$, so $DC$ (is) to $CF$. And (the latter) is converted (from the former).]\\

\begin{center}
{\large Corollary}$^\ddag$
\end{center}\vspace*{-7pt}

So (it is) clear, from this, that if composed magnitudes are proportional,
 (then) they will also be proportional (when) converted.
(Which is) the very thing it was required to show.}
\end{Parallel}
{\footnotesize \noindent$^\dag$ In modern notation, this proposition
reads that if $\alpha:\beta::\gamma:\delta$ then
$\alpha:\beta::\alpha-\gamma:\beta-\delta$.\\[0.5ex]
$^\ddag$ In modern notation, this corollary reads
that if $\alpha:\beta::\gamma:\delta$ then $\alpha:\alpha-\beta::\gamma:\gamma-\delta$.}

%%%%%%
% Prop 5.20
%%%%%%
\pdfbookmark[1]{Proposition 5.20}{pdf5.20}
\begin{Parallel}{}{} 
\ParallelLText{
\begin{center}
{\large \ggn{20}.}
\end{center}\vspace*{-7pt}

{\greekfont Ἐὰν ᾖ τρία μεγέθη καὶ ἄλλα α ὐτοῖς ἴσα τὸ
πλῆθος, σύνδυο λαμβανόμενα καὶ ἐν τῷ α ὐτῷ
λόγω, δι ἴσου δὲ τὸ πρῶτον τοῦ τρίτου μεῖζον
ᾖ, καὶ τὸ τέταρτον τοῦ ἕκτου μεῖζον ἔσται,
κἂν ἴσον, ἴσον, κἂν ἔλαττον, ἔλαττον.}\\~\\~\\

\epsfysize=0.7in
\centerline{\epsffile{Book05/fig20g.eps}}

{\greekfont Ἔστω τρία μεγέθη τὰ Α, Β, Γ, καὶ ἄλλα α ὐτοῖς
ἴσα τὸ πλῆθος τὰ Δ, Ε, Ζ, σύνδυο λαμβανόμενα
ἐν τῷ α ὐτῷ λόγῳ, ὡς μὲν τὸ Α πρὸς τὸ Β,
οὕτως τὸ Δ πρὸς τὸ Ε, ὡς δὲ τὸ Β πρὸς τὸ Γ,
οὕτως τὸ Ε πρὸς τὸ Ζ, δι ἴσου δὲ μεῖζον ἔστω
τὸ Α τοῦ Γ;
 λέγω,
ὅτι καὶ τὸ Δ τοῦ Ζ μεῖζον ἔσται, κἂν ἴσον, ἴσον,
κἂν ἔλαττον, ἔλαττον.}

{\greekfont Ἐπεὶ γὰρ μεῖζόν ἐστι τὸ Α τοῦ Γ, ἄλλο δέ τι τὸ
Β, τὸ δὲ μεῖζον πρὸς τὸ α ὐτὸ μείζονα λόγον
ἔχει ἤπερ τὸ ἔλαττον, τὸ Α ἄρα πρὸς τὸ Β μείζονα λόγον
ἔχει ἤπερ τὸ Γ πρὸς τὸ Β. ἀλλ ὡς μὲν τὸ Α πρὸς τὸ Β
[οὕτως] τὸ Δ πρὸς τὸ Ε, ὡς δὲ τὸ Γ πρὸς τὸ Β, ἀνάπαλιν
οὕτως τὸ Ζ πρὸς τὸ Ε; καὶ τὸ Δ ἄρα πρὸς τὸ Ε μείζονα λόγον
ἔχει ἤπερ τὸ Ζ πρὸς τὸ Ε. τῶν δὲ πρὸς τὸ α ὐτὸ λόγον ἐθόντων
τὸ μείζονα λόγον ἔθον μεῖζόν ἐστιν. μεῖζον ἄρα τὸ Δ τοῦ Ζ.
ὁμοίως δὴ δείχομεν, ὅτι κἂν ἴσον ᾖ τὸ Α τῷ Γ, ἴσον
ἔσται καὶ τὸ Δ τῷ Ζ, κἂν ἔλαττον, ἔλαττον.}

{\greekfont Ἐὰν ἄρα ᾖ τρία μεγέθη καὶ ἄλλα α ὐτοῖς ἴσα τὸ
πλῆθος, σύνδυο λαμβανόμενα καὶ ἐν τῷ α ὐτῷ
λόγω, δι ἴσου δὲ τὸ πρῶτον τοῦ τρίτου μεῖζον
ᾖ, καὶ τὸ τέταρτον τοῦ ἕκτου μεῖζον ἔσται,
κἂν ἴσον, ἴσον, κἂν ἔλαττον, ἔλαττον; ὅπερ ἔδει
δεῖχαι.}}

\ParallelRText{
\begin{center}
{\large Proposition 20}$^\dag$
\end{center}

If there are three magnitudes, and others of equal
 number to them,   (being) also in the same ratio taken two by two, and (if),
via equality,  the first is greater
than the third (then) the fourth will also be greater than the sixth. And
 if (the first is) equal (to the third then the fourth will also be) equal
(to the sixth). 
And if (the first is) less (than the third then the fourth will also be) less (than
the sixth).

\epsfysize=0.7in
\centerline{\epsffile{Book05/fig20e.eps}}

Let $A$, $B$, and $C$ be three magnitudes, and $D$, $E$, $F$ other (magnitudes) of equal
number to them, (being) in the same ratio taken two by two, (so that)
as $A$ (is) to $B$, so $D$ (is) to $E$, and as $B$ (is) to $C$, so $E$ (is) to $F$. And let $A$
be greater than $C$, via equality. I say that $D$ will also be greater
than $F$. And if ($A$ is) equal (to $C$ then $D$ will also be) equal (to $F$). And
if ($A$ is) less (than $C$ then $D$ will also be) less (than $F$).

For since $A$ is greater than $C$, and $B$ some other (magnitude), and
the greater (magnitude) has a greater ratio than the lesser to the same (magnitude) [Prop.~5.8], $A$ thus has
a greater ratio to $B$ than $C$ (has) to $B$. But as $A$ (is) to $B$, [so] $D$ (is)
to $E$. And, inversely, as $C$ (is) to $B$,  so  $F$ (is) to $E$ [Prop.~5.7~corr.]. Thus, $D$ also has a greater
ratio to $E$ than $F$ (has) to $E$ [Prop.~5.13]. And for (magnitudes) having a ratio to the
same (magnitude), that having the greater ratio is greater [Prop.~5.10]. Thus, $D$ (is) greater than $F$.
Similarly, we can show that even if $A$ is equal to $C$ (then) $D$ will also
be equal to $F$, and even if ($A$ is) less (than $C$ then $D$ will also
be) less (than $F$).

Thus, if there are three magnitudes, and others of equal
number to them,   (being) also in the same ratio taken two by two, and (if),
via equality, the first is greater
than the third, (then) the fourth will also be greater than the sixth. And
 if (the first is) equal (to the third then the fourth will also be) equal
(to the sixth). And (if the first is) less (than the third then the fourth will also be) less (than
the sixth). (Which is) the very thing it was required to show.}
\end{Parallel}
{\footnotesize \noindent$^\dag$ In modern notation, this proposition
reads that if $\alpha:\beta::\delta:\epsilon$ and $\beta:\gamma::\epsilon:\zeta$ then $\alpha\gtreqqless\gamma$ as $\delta\gtreqqless\zeta$.}

%%%%%%
% Prop 5.21
%%%%%%
\pdfbookmark[1]{Proposition 5.21}{pdf5.21}
\begin{Parallel}{}{} 
\ParallelLText{
\begin{center}
{\large \ggn{21}.}
\end{center}\vspace*{-7pt}

{\greekfont Ἐὰν ᾖ τρία μεγέθη καὶ ἄλλα α ὐτοῖς ἴσα τὸ πλῆθος σύνδυο λαμβανόμενα καὶ ἐν τῷ α ὐτῷ λόγῳ, ᾖ δὲ τεταραγμένη α ὐτῶν
ἡ ἀναλογία, δι ἴσου δὲ τὸ πρῶτον τοῦ τρίτου μεῖζον ᾖ, καὶ
τὸ τέταρτον τοῦ ἕκτου μεῖζον ἔσται, κἂν ἴσον, ἴσον, κἂν
ἔλαττον, ἔλαττον.}\\~\\~\\

\epsfysize=0.7in
\centerline{\epsffile{Book05/fig21g.eps}}

{\greekfont Ἔστω τρία μεγέθη τὰ Α, Β, Γ καὶ ἄλλα α ὐτοῖς ἴσα τὸ πλῆθος
τὰ Δ, Ε, Ζ, σύνδυο λαμβανόμενα καὶ ἐν τῷ α ὐτῷ λόγῳ, ἔστω δὲ
τεταραγμένη α ὐτῶν ἡ ἀναλογία, ὡς μὲν τὸ Α πρὸς τὸ Β, οὕτως
τὸ Ε πρὸς τὸ Ζ, ὡς δὲ τὸ Β πρὸς τὸ Γ, οὕτως τὸ Δ πρὸς τὸ Ε,
δι ἴσου δὲ τὸ Α τοῦ Γ μεῖζον ἔστω; λέγω, ὅτι
καὶ τὸ Δ τοῦ Ζ μεῖζον ἔσται, κἂν ἴσον, ἴσον, κἂν ἒλαττον,
ἒλαττον.}

{\greekfont Ἐπεὶ γὰρ μεῖζόν ἐστι τὸ Α τοῦ Γ, ἄλλο δέ τι τὸ Β, τὸ Α ἄρα
πρὸς τὸ Β μείζονα λόγον ἔχει ἤπερ τὸ Γ πρὸς τὸ Β. ἀλλ ὡς
μὲν τὸ Α πρὸς τὸ Β, οὕτως τὸ Ε πρὸς τὸ Ζ, ὡς δὲ τὸ Γ πρὸς
τὸ Β, ἀνάπαλιν οὕτως τὸ Ε πρὸς τὸ Δ. καὶ τὸ Ε ἄρα πρὸς τὸ Ζ
μείζονα λόγον ἔχει ἤπερ τὸ Ε πρὸς τὸ Δ. πρὸς ὃ δὲ τὸ α ὐτὸ μείζονα λόγον ἔχει, ἐκεῖνο ἔλασσόν ἐστιν; ἔλασσον ἄρα
ἐστὶ τὸ Ζ τοῦ Δ; μεῖζον ἄρα ἐστὶ τὸ Δ τοῦ Ζ.
ὁμοίως δὴ δείχομεν, ὅτι κἂν ἴσον ᾖ
τὸ Α τῷ Γ, ἴσον ἔσται καὶ τὸ Δ τῷ Ζ, κἂν ἔλαττον, ἔλαττον.}

{\greekfont Ἐὰν ἄρα ᾖ τρία μεγέθη καὶ ἄλλα α ὐτοῖς ἴσα τὸ πλῆθος, σύνδυο λαμβανόμενα καὶ ἐν τῷ α ὐτῷ λόγῳ, ᾖ δὲ τεταραγμένη α ὐτῶν
ἡ ἀναλογία, δι ἴσου δὲ τὸ πρῶτον τοῦ τρίτου μεῖζον ᾖ, καὶ
τὸ τέταρτον τοῦ ἕκτου μεῖζον ἔσται, κἂν ἴσον, ἴσον, κἂν
ἔλαττον, ἔλαττον; ὅπερ ἔδει δεῖχαι.}}

\ParallelRText{
\begin{center}
{\large Proposition 21}$^\dag$
\end{center}

If there are three magnitudes, and others of equal
 number to them, (being) also in the same ratio taken two by two, and (if)
their proportion (is) perturbed, and (if), via equality,
the first is greater than the third (then) the fourth will also
be greater than the sixth. And if (the first is) equal (to the third then the
fourth will also be) equal (to the sixth).
And if (the first is) less (than the third then the fourth will also be) less (than
the sixth).

\epsfysize=0.7in
\centerline{\epsffile{Book05/fig21e.eps}}

Let $A$, $B$, and $C$ be three magnitudes, and $D$, $E$, $F$ other (magnitudes)
of equal number to them, (being) in the same ratio taken two by two. And let their proportion be perturbed, (so that) as $A$ (is) to $B$, so $E$ (is) to $F$, and as
$B$ (is) to $C$, so $D$ (is) to $E$. And let $A$  be greater than $C$, via equality. I
say that $D$ will also be greater than $F$.
 And if ($A$ is) equal (to $C$ then $D$ will also be) equal (to $F$). And
if ($A$ is) less (than $C$ then $D$ will also be) less (than $F$).

For since $A$ is greater than $C$, and $B$ some other (magnitude),  $A$ thus
has a greater ratio to $B$ than $C$ (has) to $B$ [Prop.~5.8]. But as $A$ (is) to $B$,
so $E$ (is) to $F$. And, inversely, as $C$ (is) to $B$, so $E$ (is) to $D$ [Prop.~5.7~corr.]. Thus, $E$ also has a greater
ratio to $F$ than $E$ (has) to $D$ [Prop.~5.13]. And that (magnitude) to which the same (magnitude)  has a greater ratio  is (the) lesser (magnitude) [Prop.~5.10].
Thus, $F$ is less than $D$. Thus, $D$ is greater than $F$. Similarly, we can show that
even if $A$ is equal to $C$ (then) $D$ will also
be equal to $F$, and even if ($A$ is) less (than $C$ then $D$ will also
be) less (than $F$).

Thus, if there are three magnitudes, and others of equal
 number to them, (being) also in the same ratio taken two by two, and (if)
their proportion (is) perturbed, and (if), via equality,
the first is greater than the third (then) the fourth will also
be greater than the sixth. And if (the first is) equal (to the third then the fourth
will also be) equal (to the sixth).
And if (the first is) less (than the third then the fourth will also be) less (than
the sixth). (Which is) the very thing it was required to show.}
\end{Parallel}
{\footnotesize \noindent$^\dag$ In modern notation, this proposition
reads that if $\alpha:\beta::\epsilon:\zeta$ and $\beta:\gamma::\delta:\epsilon$ then $\alpha\gtreqqless\gamma$ as $\delta\gtreqqless\zeta$.}

%%%%%%
% Prop 5.22
%%%%%%
\pdfbookmark[1]{Proposition 5.22}{pdf5.22}
\begin{Parallel}{}{} 
\ParallelLText{
\begin{center}
{\large \ggn{22}.}
\end{center}\vspace*{-7pt}

{\greekfont Ἐὰν ᾖ ὁποσαοῦν μεγέθη καὶ ἄλλα α ὐτοῖς ἴσα τὸ πλῆθος,
σύνδυο λαμβανόμενα καὶ ἐν τῷ α ὐτῷ λόγῳ, καὶ δι ἴσου
ἐν τῷ α ὐτῷ λόγῳ ἔσται.}\\

\epsfysize=0.58in
\centerline{\epsffile{Book05/fig22g.eps}}

{\greekfont Ἔστω ὁποσαοῦν μεγέθη τὰ Α, Β, Γ καὶ ἄλλα α ὐτοῖς ἴσα τὸ
πλῆθος τὰ Δ, Ε, Ζ, σύνδυο λαμβανόμενα ἐν τῷ α ὐτῷ λόγῳ,
ὡς μὲν τὸ Α πρὸς τὸ Β, οὕτως τὸ Δ πρὸς τὸ Ε, ὡς δὲ
τὸ Β πρὸς τὸ Γ, οὕτως τὸ Ε πρὸς τὸ Ζ; λέγω, ὅτι
καὶ δι ἴσου ἐν τῷ α ὐτῳ λόγῳ
ἔσται.}

{\greekfont Εἰλήφθω γὰρ τῶν μὲν Α, Δ ἰσάκις πολλαπλάσια
τὰ Η, Θ, τῶν δὲ Β, Ε ἄλλα, ἃ ἔτυθεν, ἰσάκις πολλαπλάσια
τὰ Κ, Λ, καὶ ἔτι τῶν Γ, Ζ ἄλλα, ἃ ἔτυθεν, ἰσάκις πολλαπλάσια
τὰ Μ, Ν.}

{\greekfont Καὶ ἐπεί  ἐστιν ὡς τὸ Α πρὸς τὸ Β, οὕτως τὸ Δ πρὸς τὸ Ε,
καὶ εἴληπται τῶν μὲν Α, Δ ἰσάκις πολλαπλάσια τὰ
Η, Θ, τῶν δὲ Β, Ε ἄλλα, ἃ ἔτυθεν, ἰσάκις πολλαπλάσια
τὰ Κ, Λ, ἔστιν ἄρα ὡς τὸ Η πρὸς τὸ Κ, οὕτως τὸ Θ πρὸς τὸ Λ.
δὶα τὰ α ὐτὰ δὴ καὶ ὡς τὸ Κ πρὸς τὸ Μ, οὕτως τὸ
Λ πρὸς
τὸ Ν. ἐπεὶ οὖν τρία μεγέθη ἐστὶ τὰ Η, Κ, Μ, καὶ ἄλλα
α ὐτοῖς ἴσα
τὸ πλῆθος τὰ Θ, Λ, Ν, σύνδυο λαμβανόμενα καὶ ἐν τῷ α ὐτῷ
λόγῳ, δι ἴσου ἄρα, εἰ ὑπερέθει τὸ Η τοῦ Μ,
ὑπερέθει καὶ τὸ Θ τοῦ Ν, καὶ εἰ ἴσον, ἴσον, καὶ
εἰ ἔλαττον, ἔλαττον. καί ἐστι τὰ μὲν Η, Θ τῶν Α, Δ
ἰσάκις πολλαπλάσια, τὰ δὲ Μ, Ν τῶν Γ, Ζ ἄλλα, ἃ
ἔτυθεν, ἰσάκις πολλαπλάσια. ἔστιν ἄρα ὡς τὸ Α πρὸς
τὸ Γ, οὕτως τὸ Δ πρὸς τὸ Ζ.}

{\greekfont Ἐὰν ἄρα ᾖ ὁποσαοῦν μεγέθη καὶ ἄλλα α ὐτοῖς ἴσα τὸ πλῆθος,
σύνδυο λαμβανόμενα ἐν τῷ α ὐτῷ λόγῳ, καὶ δι ἴσου
ἐν τῷ α ὐτῷ λόγῳ ἔσται; ὅπερ ἔδει δεῖχαι.}}

\ParallelRText{
\begin{center}
{\large Proposition 22}$^\dag$
\end{center}

If there are any number of magnitudes
whatsoever, and (some) other (magnitudes)  of equal  number to them, (which are)
also in the same ratio taken two by two, (then) they will also be in the
same ratio via equality.

\epsfysize=0.58in
\centerline{\epsffile{Book05/fig22e.eps}}

Let there be any number of magnitudes whatsoever, $A$, $B$,  $C$, and
(some) other (magnitudes), $D$, $E$, $F$, of equal number to them, (which are)
in the same ratio taken two by two, (so that) as $A$ (is) to $B$, so
$D$ (is) to $E$, and as $B$ (is) to $C$, so $E$ (is) to $F$. I say that they
will also be in the same ratio via equality. (That is, as $A$ is to $C$, so
$D$ is to $F$.)

For let the equal multiples $G$ and $H$ be taken of $A$ and $D$ (respectively), and the other random equal multiples $K$ and $L$ of $B$ and
$E$ (respectively), and the yet other random equal multiples $M$ and $N$ of
$C$ and $F$ (respectively).

And since as $A$ is to $B$, so $D$ (is) to $E$, and the equal multiples $G$ and $H$
have been taken of $A$ and $D$ (respectively), and the other random
equal multiples $K$ and $L$ of $B$ and $E$ (respectively), thus as $G$ is to $K$,
so $H$ (is) to $L$ [Prop.~5.4]. And, so, for the same
(reasons), as $K$ (is) to $M$, so $L$ (is) to $N$. Therefore, since $G$, $K$, and $M$ are three magnitudes, and $H$, $L$, and $N$ other (magnitudes) of equal number to them,
(which are) also in the same ratio taken two by two, thus, via equality,
if $G$ exceeds $M$ (then) $H$ also exceeds $N$, and if ($G$ is) equal (to $M$ then $H$
is also) equal (to $N$), and if ($G$ is) less (than $M$ then $H$
is also) less (than $N$) [Prop.~5.20].
And  $G$ and $H$ are equal multiples of $A$ and $D$ (respectively),
and $M$ and $N$ other random equal multiples of $C$ and $F$ (respectively).
Thus, as $A$ is to $C$, so $D$ (is) to $F$ [Def.~5.5].

Thus, if there are any number of magnitudes
whatsoever, and (some) other (magnitudes) of equal number to them, (which are)
also in the same ratio taken two by two, (then) they will also be in the
same ratio via equality. (Which is) the very thing
it was required to show.}
\end{Parallel}
{\footnotesize \noindent $^\dag$ In modern notation, this proposition
reads that if $\alpha:\beta::\epsilon:\zeta$ and $\beta:\gamma::\zeta:\eta$ 
and $\gamma:\delta::\eta:\theta$ 
then $\alpha:\delta::\epsilon:\theta$.}

%%%%%%
% Prop 5.23
%%%%%%
\pdfbookmark[1]{Proposition 5.23}{pdf5.23}
\begin{Parallel}{}{} 
\ParallelLText{
\begin{center}
{\large\ggn{23}.}
\end{center}\vspace*{-7pt}

{\greekfont Ἐὰν ᾖ τρία μεγέθη καὶ ἄλλα α ὐτοῖς ἴσα τὸ πλῆθος σύνδυο
λαμβανόμενα ἐν τῷ α ὐτῷ λόγῳ, ᾖ δὲ τεταραγμένη α ὐτῶν
ἡ ἀναλογία, καὶ δι ἴσου ἐν τῷ α ὐτῷ λόγῳ ἔσται.}\\

\epsfysize=0.75in
\centerline{\epsffile{Book05/fig23g.eps}}

{\greekfont Ἔστω τρία μεγέθη τὰ Α, Β, Γ καὶ ἄλλα α ὐτοῖς ἴσα τὸ πλῆθος 
σύνδυο λαμβανόμενα ἐν τῷ α ὐτῷ λόγῳ τὰ Δ, Ε, Ζ, ἔστω
δὲ τεταραγμένη α ὐτῶν ἡ ἀναλογία, ὡς μὲν τὸ Α πρὸς τὸ Β,
οὕτως τὸ Ε πρὸς τὸ Ζ, ὡς δὲ τὸ Β πρὸς τὸ Γ, οὕτως τὸ Δ
πρὸς τὸ Ε; λέγω, ὅτι ἐστὶν ὡς τὸ Α πρὸς τὸ Γ, οὕτως
τὸ Δ πρὸς τὸ Ζ.}

{\greekfont Εἰλήφθω τῶν μὲν Α, Β, Δ ἰσάκις πολλαπλάσια τὰ Η, Θ, Κ, τῶν
δὲ Γ, Ε, Ζ ἄλλα, ἃ ἔτυθεν, ἰσάκις πολλαπλάσια τὰ Λ, Μ, Ν.}

{\greekfont Καὶ ἐπεὶ ἰσάκις ἐστὶ πολλαπλάσια τὰ Η, Θ τῶν Α, Β, τὰ δὲ μέρη
τοὶς ὡσαύτως πολλαπλασίοις τὸν α ὐτὸν ἔχει λόγον,
ἔστιν ἄρα ὡς τὸ Α πρὸς τὸ Β, οὕτως τὸ Η πρὸς
τὸ Θ. διὰ τὰ α ὐτὰ δὴ καὶ ὡς τὸ Ε πρὸς τὸ Ζ, οὕτως τὸ Μ
πρὸς τὸ Ν; καί ἐστιν ὡς τὸ Α πρὸς τὸ Β, οὕτως τὸ Ε πρὸς
τὸ Ζ; καὶ ὡς ἄρα τὸ Η πρὸς τὸ Θ, οὕτως τὸ Μ πρὸς τὸ Ν.
καὶ ἐπεί ἐστιν ὡς τὸ Β πρὸς τὸ Γ, οὕτως τὸ Δ πρὸς τὸ Ε,
καὶ ἐναλλὰχ ὡς τὸ Β πρὸς τὸ Δ, οὕτως τὸ Γ πρὸς τὸ Ε. καὶ
ἐπεὶ τὰ Θ, Κ τῶν Β, Δ ἰσάκις ἐστὶ πολλαπλάσια, τὰ δὲ μέρη 
τοῖς ἰσάκις πολλαπλασίοις τὸν α ὐτὸν ἔχει λόγον, ἔστιν ἄρα
ὡς τὸ Β πρὸς τὸ Δ, οὕτως τὸ Θ πρὸς τὸ Κ. ἀλλ ὡς 
τὸ Β πρὸς τὸ Δ, οὕτως τὸ Γ πρὸς τὸ Ε; καὶ ὡς ἄρα τὸ Θ πρὸς
τὸ Κ, οὕτως τὸ Γ πρὸς τὸ Ε.
 πάλιν, ἐπεὶ τὰ Λ, Μ
τῶν Γ, Ε ἰσάκις ἐστι πολλαπλάσια, ἔστιν ἄρα ὡς τὸ Γ
πρὸς τὸ Ε, οὕτως τὸ Λ πρὸς τὸ Μ. ἀλλ ὡς τὸ Γ πρὸς τὸ Ε,
οὕτως τὸ Θ πρὸς τὸ Κ; καὶ ὡς ἄρα τὸ Θ πρὸς τὸ Κ, οὕτως
τὸ Λ πρὸς τὸ Μ, καὶ ἐναλλὰχ ὡς τὸ Θ πρὸς τὸ Λ, τὸ Κ
πρὸς τὸ Μ. ἐδείθθη δὲ καὶ ὡς τὸ Η πρὸς τὸ Θ, οὕτως
τὸ Μ πρὸς τὸ Ν. ἐπεὶ οὖν τρία μεγέθη ἐστὶ τὰ Η, Θ, Λ,
καὶ ἄλλα α ὐτοις ἴσα τὸ πλῆθος τὰ Κ, Μ, Ν σύνδυο λαμβανόμενα
ἐν τῷ α ὐτῷ λόγῳ, καί ἐστιν α ὐτῶν τεταραγμένη ἡ
ἀναλογία, δι ἴσου ἄρα, εἰ ὑπερέθει τὸ Η
τοῦ Λ, ὑπερέθει καὶ τὸ Κ τοῦ Ν, καὶ εἰ ἴσον, ἴσον, καὶ
εἰ ἔλαττον, ἔλαττον. καί ἐστι τὰ μὲν Η, Κ τῶν Α, Δ ἰσάκις
πολλαπλάσια, τὰ δὲ Λ, Ν τῶν Γ, Ζ. ἔστιν ἄρα ὡς τὸ Α πρὸς
τὸ Γ, οὕτως τὸ Δ πρὸς τὸ Ζ.}

{\greekfont Ἐὰν ἄρα  ᾖ τρία μεγέθη καὶ ἄλλα α ὐτοῖς ἴσα τὸ πλῆθος σύνδυο
λαμβανόμενα ἐν τῷ α ὐτῷ λόγῳ, ᾖ δὲ τεταραγμένη α ὐτῶν
ἡ ἀναλογία, καὶ δι ἴσου ἐν τῷ α ὐτῷ λόγῳ ἔσται;
ὅπερ ἔδει δεῖχαι.}}

\ParallelRText{
\begin{center}
{\large Proposition 23}$^\dag$
\end{center}

If there are three magnitudes, and others of equal number to them, (being) in the same ratio taken two by two, and (if) their
proportion is perturbed, (then) they will also be in the same ratio via equality.

\epsfysize=0.75in
\centerline{\epsffile{Book05/fig23e.eps}}

Let $A$, $B$, and $C$ be three magnitudes, and $D$, $E$ and $F$ other (magnitudes) of equal  number to them, (being) in the same ratio taken two by two.
And let their proportion be perturbed, (so that) as $A$ (is) to $B$, so $E$ (is) to $F$,
and as $B$ (is) to $C$, so $D$ (is) to $E$. I say that as $A$ is to $C$, so $D$ (is) to $F$.

Let the equal multiples $G$, $H$, and $K$ be taken of $A$, $B$, and $D$ (respectively), and the other random equal multiples $L$, $M$, and $N$ of
$C$, $E$, and $F$ (respectively).

And since $G$ and $H$ are equal multiples of $A$ and $B$ (respectively), and
parts have the same ratio as similar multiples [Prop.~5.15], thus as $A$ (is) to $B$, so $G$ (is) to $H$. 
And, so, for the same (reasons), as $E$ (is) to $F$, so $M$ (is) to $N$. And as
$A$ is to $B$, so $E$ (is) to $F$.  And, thus, as $G$ (is) to $H$, so $M$ (is) to $N$ [Prop.~5.11]. And since as $B$ is to $C$, so $D$ (is) to $E$,
also, alternately, as $B$ (is) to $D$, so $C$ (is) to $E$ [Prop.~5.16]. And since $H$ and $K$ are equal multiples
of $B$ and $D$ (respectively), and parts have the same ratio as similar multiples
[Prop.~5.15], thus as $B$ is to $D$, so $H$
(is) to $K$. But, as $B$ (is) to $D$, so $C$ (is) to $E$.  And, thus, as $H$ (is) to $K$, so $C$ (is) to $E$ [Prop.~5.11]. Again, since
$L$ and $M$ are  equal multiples of $C$ and $E$ (respectively), thus as
$C$ is to $E$, so $L$ (is) to $M$ [Prop.~5.15].
But, as $C$ (is) to $E$, so $H$ (is) to $K$. And, thus, as $H$ (is) to $K$,
so $L$ (is) to $M$ [Prop.~5.11]. Also, alternately,
as $H$ (is) to $L$, so $K$ (is) to $M$ [Prop.~5.16].
And it was also shown (that) as $G$ (is) to $H$, so $M$ (is) to $N$. Therefore,
since $G$, $H$, and $L$ are three magnitudes, and $K$, $M$, and $N$ other (magnitudes)
of equal number to them, (being) in the same ratio taken two by two, 
and their proportion is perturbed, thus, via equality, if $G$ exceeds
$L$ (then) $K$ also exceeds $N$, and if ($G$ is) equal (to $L$ then $K$ is also) equal (to $N$),
and if ($G$ is) less (than $L$ then $K$ is also) less (than $N$)  [Prop.~5.21]. And $G$ and $K$ are equal
multiples of $A$ and $D$ (respectively), and $L$ and $N$ of $C$ and $F$ (respectively). 
Thus, as $A$ (is) to $C$, so $D$ (is) to $F$  [Def.~5.5].

Thus, if there are three magnitudes, and others of  equal number to them, (being) in the same ratio taken two by two, and (if) their
proportion is perturbed, (then) they will also be in the same ratio via equality.
(Which is) the very thing it was required to show.}
\end{Parallel}
{\footnotesize \noindent$^\dag$ In modern notation, this proposition
reads that if $\alpha:\beta::\epsilon:\zeta$ and $\beta:\gamma::\delta:\epsilon$ then $\alpha:\gamma::\delta:\zeta$.}

%%%%%%
% Prop 5.24
%%%%%%
\pdfbookmark[1]{Proposition 5.24}{pdf5.24}
\begin{Parallel}{}{} 
\ParallelLText{
\begin{center}
{\large \ggn{24}.}
\end{center}\vspace*{-7pt}

{\greekfont Ἐὰν πρῶτον πρὸς δεύτερον τὸν α ὐτὸν ἔθῃ λόγον καὶ τρίτον
πρὸς τέταρτον, ἔθῃ δὲ καὶ πέμπτον πρὸς δεύτερον τὸν α ὐτὸν
λόγον καὶ ἕκτον πρὸς τέταρτον, καὶ συντεθὲν πρῶτον καὶ
πέμπτον πρὸς δεύτερον τὸν α ὐτὸν ἕχει λόγον καὶ τρίτον
καὶ ἕκτον πρὸς τέταρτον.}

{\greekfont Πρῶτον γὰρ τὸ ΑΒ πρὸς δεύρερον τὸ Γ τὸν α ὐτὸν ἐθέτω
λόγον καὶ τρίτον τὸ ΔΕ πρὸς τέταρτον τὸ Ζ, ἐθέτω δὲ καὶ πέμπτον
τὸ ΒΗ πρὸς δεύτερον τὸ Γ τὸν α ὐτὸν λόγον καὶ ἕκτον τὸ ΕΘ
πρὸς τέταρτον τὸ Ζ; λέγω, ὅτι καὶ συντεθὲν πρῶτον καὶ πέμπτον
τὸ ΑΗ πρὸς δεύτερον τὸ Γ τὸν α ὐτὸν ἕχει λόγον, καὶ
τρίτον καὶ ἕκτον τὸ ΔΘ πρὸς τέταρτον τὸ Ζ.}\\~\\

\epsfysize=1.4in
\centerline{\epsffile{Book05/fig24g.eps}}

{\greekfont Ἐπεὶ γάρ ἐστιν ὡς τὸ ΒΗ πρὸς τὸ Γ, οὕτως τὸ ΕΘ πρὸς τὸ Ζ,
ἀνάπαλιν ἄρα ὡς τὸ Γ πρὸς τὸ ΒΗ, οὕτως
τὸ Ζ πρὸς τὸ ΕΘ. ἐπεὶ οὖν ἐστιν ὡς τὸ ΑΒ πρὸς τὸ Γ,
οὕτως τὸ ΔΕ πρὸς τὸ Ζ, ὡς δὲ τὸ Γ πρὸς τὸ ΒΗ, οὕτως
τὸ Ζ πρὸς τὸ ΕΘ, δι ἴσου ἄρα ἐστὶν ὡς τὸ ΑΒ πρὸς τὸ ΒΗ, οὕτως
τὸ ΔΕ πρὸς τὸ ΕΘ. καὶ ἐπεὶ διῃρημένα μεγέθη ἀνάλογόν ἐστιν,
καὶ συντεθέντα ἀνάλογον ἔσται; ἔστιν ἄρα ὡς τὸ ΑΗ πρὸς τὸ ΗΒ,
οὕτως τὸ ΔΘ πρὸς τὸ ΘΕ. ἔστι δὲ καὶ ὡς τὸ ΒΗ πρὸς τὸ Γ, οὕτως
τὸ ΕΘ πρὸς τὸ Ζ;
δι ἴσου ἄρα ἐστὶν ὡς τὸ ΑΗ πρὸς τὸ Γ, οὕτως τὸ
ΔΘ πρὸς τὸ Ζ.}

{\greekfont Ἐὰν ἄρα πρῶτον πρὸς δεύτερον τὸν α ὐτὸν ἔθῃ λόγον καὶ τρίτον
πρὸς τέταρτον, ἔθῃ δὲ καὶ πέμπτον πρὸς δεύτερον τὸν α ὐτὸν
λόγον καὶ ἕκτον πρὸς τέταρτον, καὶ συντεθὲν πρῶτον καὶ
πέμπτον πρὸς δεύτερον τὸν α ὐτὸν ἕχει λόγον καὶ τρίτον
καὶ ἕκτον πρὸς τέταρτον; ὅπερ ἔδει δεὶχαι.}}

\ParallelRText{
\begin{center}
{\large Proposition 24}$^\dag$
\end{center}

If a first (magnitude) has to a second the same
ratio  that third (has) to a fourth, and a fifth (magnitude)
also has to the second the same ratio that a sixth  (has) to the fourth, (then)
  the first (magnitude) and the fifth, added together, will also have the same ratio to the
 second
 that the third (magnitude) and sixth  (added together, have) to the fourth.
 
 For let a first (magnitude) $AB$ have the same ratio to a second $C$
 that a third $DE$ (has) to a fourth $F$. And let a fifth (magnitude) $BG$
 also have the same ratio to the second $C$ that a sixth $EH$ (has) to
 the fourth $F$. I say that the first (magnitude) and the fifth, added together, $AG$, will also  have the same ratio to the second $C$ that the third (magnitude) and the sixth, (added together), $DH$,  (has) to the fourth $F$.
 
\epsfysize=1.4in
\centerline{\epsffile{Book05/fig24e.eps}}

 For since as $BG$ is to $C$, so $EH$ (is) to $F$,  thus, inversely, as $C$ (is) to $BG$,
 so $F$ (is) to $EH$  [Prop.~5.7~corr.]. Therefore,
 since as $AB$ is to $C$, so $DE$ (is) to $F$, and as $C$ (is) to $BG$, so
 $F$ (is) to $EH$, thus, via equality, as $AB$ is to $BG$, so $DE$ (is) to $EH$ 
[Prop.~5.22]. And since separated magnitudes
are proportional, (then) they will also be proportional (when) composed [Prop.~5.18]. Thus, as $AG$ is to $GB$, so $DH$
(is) to $HE$. And, also, as $BG$ is to $C$, so $EH$ (is) to $F$. Thus, via equality, 
as $AG$ is to $C$, so $DH$ (is) to $F$ [Prop.~5.22].

Thus, if a first (magnitude) has to a second the same
ratio that a third (has) to a fourth, and a fifth (magnitude)
also has to the second the same ratio that a sixth  (has) to the fourth, (then)
the first (magnitude) and the fifth, added together, will also have the same ratio to the
 second
 that the third (magnitude) and the sixth  (added together, have) to the fourth. (Which is) the very thing it was required to show.}
\end{Parallel}
{\footnotesize \noindent$^\dag$ In modern notation, this proposition
reads that if $\alpha:\beta::\gamma:\delta$ and $\epsilon:\beta::\zeta:\delta$ then $\alpha+\epsilon:\beta::\gamma+\zeta:\delta$.}

%%%%%%
% Prop 5.25
%%%%%%
\pdfbookmark[1]{Proposition 5.25}{pdf5.25}
\begin{Parallel}{}{} 
\ParallelLText{
\begin{center}
{\large \ggn{25}.}
\end{center}\vspace*{-7pt}

{\greekfont Ἐὰν τέσσαρα μεγέθη ἀνάλογον ᾖ, τὸ μέγιστον [α ὐτῶν] καὶ
τὸ ἐλάθιστον δύο τῶν λοιπῶν μείζονά ἐστιν.}\\

\epsfysize=1.4in
\centerline{\epsffile{Book05/fig25g.eps}}

{\greekfont Ἔστω τέσσαρα μεγέθη ἀνάλογον τὰ ΑΒ, ΓΔ, Ε, Ζ, ὡς τὸ ΑΒ
πρὸς τὸ ΓΔ, οὕτως τὸ Ε πρὸς τὸ Ζ, ἔστω δὲ μέγιστον μὲν
α ὐτῶν τὸ ΑΒ, ἐλάθιστον δὲ τὸ Ζ; λέγω, ὅτι τὰ ΑΒ, Ζ τῶν
ΓΔ, Ε μείζονά ἐστιν.}

{\greekfont Κείσθω γὰρ τῷ μὲν Ε ἴσον τὸ ΑΗ, τῷ δὲ Ζ
ἴσον τὸ ΓΘ.}

{\greekfont Ἐπεὶ [οὖν] ἐστιν ὡς τὸ ΑΒ πρὸς τὸ ΓΔ, οὕτως
τὸ Ε πρὸς τὸ Ζ, ἴσον δὲ τὸ  μὲν Ε τῷ ΑΗ, τὸ δὲ Ζ τῷ ΓΘ,
ἔστιν ἄρα ὡς τὸ ΑΒ πρὸς τὸ ΓΔ, οὕτως
τὸ ΑΗ πρὸς τὸ ΓΘ. καὶ ἐπεί ἐστιν ὡς ὅλον τὸ ΑΒ πρὸς
ὅλον τὸ ΓΔ, οὕτως ἀφαιρεθὲν τὸ ΑΗ πρὸς ἀφαιρεθὲν τὸ ΓΘ,
καὶ λοιπὸν ἄρα τὸ ΗΒ πρὸς λοιπὸν τὸ ΘΔ ἔσται ὡς ὅλον
τὸ ΑΒ πρὸς ὅλον τὸ ΓΔ. μεῖζον δὲ τὸ ΑΒ τοῦ ΓΔ; μεῖζον
ἄρα καὶ τὸ ΗΒ τοῦ ΘΔ. καὶ ἐπεὶ ἴσον ἐστὶ τὸ μὲν ΑΗ τῷ
Ε, τὸ δὲ ΓΘ τῷ Ζ, τὰ ἄρα ΑΗ, Ζ ἴσα ἐστὶ τοῖς
ΓΘ, Ε.
καὶ [ἐπεὶ] ἐὰν [ἀνίσοις ἴσα προστεθῇ, τὰ ὅλα
ἄνισά ἐστιν, ἐὰν ἄρα] τῶν ΗΒ, ΘΔ ἀνίσων ὄντων καὶ
μείζονος τοῦ ΗΒ  τῷ μὲν ΗΒ προστεθῇ τὰ ΑΗ, Ζ, τῷ δὲ ΘΔ
προστεθῇ τὰ ΓΘ, Ε, συνάγεται τὰ ΑΒ, Ζ μείζονα τῶν ΓΔ, Ε.}

{\greekfont Ἐὰν ἄρα τέσσαρα μεγέθη ἀνάλογον ᾖ, τὸ μέγιστον α ὐτῶν καὶ
τὸ ἐλάθιστον δύο τῶν λοιπῶν μείζονά ἐστιν. ὅπερ ἔδει δεῖχαι.}}

\ParallelRText{
\begin{center}
{\large Proposition 25}$^\dag$
\end{center}

If four magnitudes are proportional, (then) the (sum of the)
largest  and the smallest [of them] is greater than the (sum of the)  remaining two (magnitudes).

\epsfysize=1.4in
\centerline{\epsffile{Book05/fig25e.eps}}

Let $AB$, $CD$, $E$, and $F$ be four proportional magnitudes, (such that) as
$AB$ (is) to $CD$, so $E$ (is) to $F$. And let $AB$ be the greatest of them, and $F$ the least. I say that $AB$ and $F$ is greater than $CD$ and $E$.

For let $AG$ be made equal to $E$, and $CH$ equal to $F$.

\mbox{[}In fact,] since as $AB$ is to $CD$, so $E$ (is) to $F$, and $E$ (is) equal to $AG$, and
$F$ to $CH$,  thus as $AB$ is to $CD$, so $AG$ (is) to $CH$. And since the whole
$AB$ is to the whole $CD$ as the (part) taken away $AG$ (is) to the
(part) taken away $CH$, thus the remainder $GB$ will also be to the
remainder $HD$ as the whole $AB$ (is) to the whole $CD$ [Prop.~5.19]. And $AB$ (is) greater than $CD$.
Thus, $GB$ (is) also greater than $HD$. And since $AG$ is equal to $E$, and $CH$
to $F$, thus $AG$ and $F$ is equal to $CH$ and $E$. And [since] if [equal (magnitudes)
are added to unequal (magnitudes then) the wholes are unequal, thus if] $AG$ and $F$ are added to $GB$, and $CH$ and
$E$ to $HD$---$GB$ and
$HD$ being unequal, and $GB$ greater---it is inferred that $AB$ and $F$ (is) greater than $CD$ and $E$.

Thus, if four magnitudes are proportional, (then) the (sum of the)
largest  and the smallest of them is greater than the (sum of the) remaining two (magnitudes). (Which is) the very thing it was required to show.}
\end{Parallel}
{\footnotesize \noindent$^\dag$ In modern notation, this proposition
reads that if $\alpha:\beta::\gamma:\delta$, and $\alpha$ is the greatest and $\delta$ the least, then $\alpha+\delta>\beta+\gamma$.}

\newpage~\thispagestyle{empty}
